%---------------------------------------------------------------------------------------

\chapter{Método de Lattice Boltzmann}\label{ch:metodo-de-lattice-boltzmann}

%---------------------------------------------------------------------------------------

% Lista de abreviaturas e de siglas
\criarsigla{DBE}{\textit{Discrete Boltzmann Equation}}
\criarsigla{ELBM}{\textit{Entropic Lattice Boltzmann Method}}
\criarsigla{EOS}{\textit{Equation of State}}
\criarsigla{HRR}{\textit{Hybrid Recursive Regularization}}
\criarsigla{LU}{\textit{Lattice Units}}
\criarsigla{MRT}{\textit{Multiple Relaxation Time}}
\criarsigla{RR}{\textit{Recursive Regularization}}
\criarsigla{TRT}{\textit{Two Relaxation Time}}
\criarsigla{WENO}{\textit{Weighted Essentially Non-Oscillatory}}

%---------------------------------------------------------------------------------------

% Lista de simbolos

%---------------------------------------------------------------------------------------

O método de Lattice Boltzmann parte de uma premissa distinta aos métodos macroscópicos usuais.
Em lugar de discretizar diretamente as equações de conservação, discretiza-se a equação de Boltzmann, e o comportamento hidrodinâmico emerge como consequência da dinâmica das funções de distribuição sobre uma rede regular de velocidades.
Essa construção apoia-se em dois pilares teóricos: a redução do espaço contínuo de velocidades a um conjunto finito de direções, obtida por quadratura de Gauss-Hermite, e a análise multiescala de Chapman-Enskog, que demonstra a recuperação das equações de Navier-Stokes no limite hidrodinâmico.

O regime compressível térmico, objeto deste trabalho, exige extensões dessa formulação básica.
A presença de ondas de choque e de variações acentuadas de temperatura torna insuficientes tanto o operador de colisão de tempo de relaxação único quanto a distribuição de equilíbrio truncada em segunda ordem, o que motiva o emprego de operadores de múltiplos tempos de relaxação, de técnicas de regularização e de uma segunda função de distribuição dedicada à energia.
O desenvolvimento a seguir reconstitui essa fundamentação teórica na extensão necessária para justificar cada decisão incorporada ao miniaplicativo, cuja tradução em código é tratada no Capítulo~\ref{ch:implementacao-computacional}.

%---------------------------------------------------------------------------------------

\section{Formulação Clássica do Método de Lattice Boltzmann}\label{sec:lbm-deducao}

A formulação hoje consolidada do LBM pode ser alcançada por dois caminhos que convergem para a mesma equação de evolução.
O percurso historicamente original partiu dos autômatos de gás de rede (LGA), nos quais o estado de cada nó é descrito por variáveis booleanas de ocupação $n_i \in \{0, 1\}$, atualizadas por
\begin{equation}\label{eq:lga}
    n_i(\mathbf{x} + \mathbf{c}_i, t + 1) = n_i(\mathbf{x}, t) + \Omega_i(n(\mathbf{x}, t))
\end{equation}
onde $\Omega_i$ representa as regras de colisão entre partículas \parencite{Chen1998}.
A substituição das variáveis booleanas por suas médias de conjunto (\textit{ensemble average}), $f_i = \langle n_i \rangle$, elimina o ruído estatístico que inviabilizava o LGA em escoamentos de interesse prático e preserva a estrutura local do algoritmo, dando origem ao LBM.
O segundo caminho, posterior e mais geral, parte da teoria cinética e chega à rede por discretização sucessiva do espaço de velocidades, do espaço e do tempo.
A convergência entre o percurso histórico e a discretização direta da equação de Boltzmann é o que confere ao método sua fundamentação cinética, e é esta última que se desenvolve a seguir.

A derivação do LBM a partir da equação de Boltzmann contínua baseia-se na discretização do espaço de velocidades e, em seguida, do espaço e do tempo.
Na ausência de forças externas, a equação cinética pode ser simplificada pelo modelo de relaxação de \textcite{Bhatnagar1954}, conhecido como operador de colisão BGK, que substitui o termo integral de colisão por uma relaxação linear em direção a um estado de equilíbrio termodinâmico local:
\begin{equation}
    \frac{\partial f}{\partial t} + \boldsymbol{\xi} \cdot \nabla f = -\frac{1}{\tau}(f - f^{eq})
\end{equation}
onde $\tau$ é o tempo de relaxação, que determina a taxa de retorno ao equilíbrio, e $f^{eq}$ corresponde à distribuição de equilíbrio local de Maxwell-Boltzmann, dada por:
\begin{equation}
    f^{eq}(\mathbf{x}, \boldsymbol{\xi}, t) = \frac{\rho}{(2\pi RT)^{D/2}} \exp\left( -\frac{(\boldsymbol{\xi} - \mathbf{u})^2}{2RT} \right)
\end{equation}
sendo $R$ a constante específica do gás, $T$ a temperatura e $D$ a dimensão do espaço.

Antes de proceder à discretização, é conveniente adimensionalizar a equação cinética pelas escalas macroscópicas do escoamento, isto é, pelo comprimento característico $L$, pela velocidade de referência $U$ e por uma massa específica de referência $\rho_0$, adotando-se ainda $\sqrt{RT_0}$ como escala das velocidades microscópicas.
Nessa forma adimensional, o tempo de relaxação comparece unicamente na combinação $\tau \sqrt{RT_0}/L$, proporcional ao número de Knudsen $Kn$, definido como a razão entre o percurso livre médio das partículas e o comprimento característico do escoamento.
O regime hidrodinâmico corresponde a $Kn \ll 1$, condição na qual as colisões são frequentes o bastante para manter a distribuição próxima do equilíbrio local, e é justamente essa proximidade que autoriza tanto o truncamento descrito a seguir quanto a análise multiescala apresentada na Seção~\ref{sec:lbm-chapman-enskog}.

Para transpor a formulação contínua ao âmbito numérico, o espaço contínuo de velocidades $\boldsymbol{\xi}$ deve ser restrito a um conjunto finito de velocidades discretas $\mathbf{c}_i$.
A redução é possível porque não é a função $f$ em si que possui significado hidrodinâmico, e sim um pequeno número de seus momentos, associados à massa específica, à quantidade de movimento e à energia.
Conforme demonstrado por \textcite{Shan2006}, essa observação conduz à projeção da função de distribuição em uma base de polinômios tensoriais ortogonais de Hermite:
\begin{equation}\label{eq:expansao-hermite}
    f(\mathbf{x}, \boldsymbol{\xi}, t) = \varpi(\boldsymbol{\xi}) \sum_{n=0}^{\infty} \frac{1}{n!} \mathbf{a}^{(n)}(\mathbf{x}, t) : \boldsymbol{\mathcal{H}}^{(n)}(\boldsymbol{\xi})
\end{equation}
onde $\varpi(\boldsymbol{\xi}) = (2\pi)^{-D/2} \exp(-|\boldsymbol{\xi}|^2/2)$ é a função peso gaussiana, $\boldsymbol{\mathcal{H}}^{(n)}$ é o polinômio tensorial de Hermite de ordem $n$, o símbolo $:$ denota a contração completa entre tensores de mesma ordem e $\mathbf{a}^{(n)} = \int f \boldsymbol{\mathcal{H}}^{(n)} d\boldsymbol{\xi}$ são os coeficientes da expansão.
A escolha dessa base não é arbitrária: como os polinômios de Hermite são ortogonais em relação ao peso gaussiano, e como a distribuição de equilíbrio é justamente uma gaussiana deslocada, os coeficientes $\mathbf{a}^{(n)}$ coincidem com os próprios momentos da distribuição.
Aplicada ao equilíbrio, a expansão fornece $a^{(0)} = \rho$, $\mathbf{a}^{(1)} = \rho \mathbf{u}$ e $\mathbf{a}^{(2)} = \rho \left[ \mathbf{u}\mathbf{u} + (T - 1)\mathbf{I} \right]$, o que evidencia que reter os três primeiros termos basta para representar a massa específica, a quantidade de movimento e o tensor de pressão.

O truncamento da \eqref{eq:expansao-hermite} em uma ordem $\mathcal{N}$ finita é o que torna a discretização possível, e sua consequência é direta: apenas os momentos até a ordem $\mathcal{N}$ são reproduzidos com exatidão.
Reter $\mathcal{N} = 2$ é suficiente para recuperar as equações de Navier-Stokes isotérmicas, enquanto a descrição consistente do transporte de energia e dos efeitos térmicos exige $\mathcal{N} \ge 3$, exigência que reaparece adiante como uma das razões da inadequação do modelo clássico ao regime compressível.

Truncada a expansão, todo momento de interesse assume a forma $\int \varpi(\boldsymbol{\xi}) P(\boldsymbol{\xi}) d\boldsymbol{\xi}$, com $P$ polinomial, integral que pode ser calculada exatamente por uma quadratura de Gauss-Hermite de ordem suficiente:
\begin{equation}\label{eq:quadratura-gauss-hermite}
    \int \varpi(\boldsymbol{\xi}) P(\boldsymbol{\xi}) d\boldsymbol{\xi} = \sum_{i=0}^{Q-1} w_i P(\mathbf{c}_i)
\end{equation}
onde os $Q$ pontos da quadratura $\mathbf{c}_i$ tornam-se as velocidades discretas da rede e os $w_i$, seus pesos.
Em uma dimensão, a regra de três pontos emprega as raízes do polinômio de Hermite de terceira ordem, $\xi \in \{-\sqrt{3},\, 0,\, +\sqrt{3}\}$, com pesos $\{1/6,\, 2/3,\, 1/6\}$, e integra exatamente polinômios de grau até cinco.
O produto tensorial dessa regra em $D$ dimensões gera $3^D$ pontos, o que corresponde exatamente aos conjuntos D2Q9 em duas dimensões e D3Q27 em três, com pesos dados pelo produto dos pesos unidimensionais.
Reescalando a unidade de velocidade de modo que os pontos de quadratura coincidam com os nós vizinhos da malha, ou seja, $\mathbf{c}_i = c\, \mathbf{v}_i$ com $\mathbf{v}_i$ de componentes em $\{-1, 0, +1\}$ e $c = \Delta x / \Delta t$, obtém-se a velocidade do som na rede $c_s = c/\sqrt{3}$ e os pesos $w_0 = (2/3)^2 = 4/9$, $w_i = (2/3)(1/6) = 1/9$ para as direções axiais e $w_i = (1/6)^2 = 1/36$ para as diagonais.
A \autoref{fig:lbm-quadratura} ilustra essa construção.
Os pesos usuais do LBM, longe de constituírem uma escolha empírica, são portanto os pesos de uma quadratura de Gauss-Hermite, e é essa origem que garante a preservação exata dos momentos hidrodinâmicos e, com ela, a conservação de massa específica, de quantidade de movimento e de energia.

\begin{figure}[H]
    \centering
    \caption{Obtenção do conjunto D2Q9 pelo produto tensorial da regra de Gauss-Hermite de três pontos, com os pesos resultantes em cada direção discreta.}
    \label{fig:lbm-quadratura}
    \fbox{\begin{minipage}{0.85\textwidth}
        \centering
        \vspace{2cm}
        [Espaço reservado para figura: construção da rede D2Q9 por produto tensorial. À esquerda, um eixo unidimensional com os três pontos da quadratura de Gauss-Hermite em $-\sqrt{3}$, $0$ e $+\sqrt{3}$, com os pesos $1/6$, $2/3$ e $1/6$ indicados sob cada ponto, e a curva da gaussiana de referência ao fundo. À direita, a malha de nove pontos resultante do produto dos dois eixos, cada ponto rotulado com o peso correspondente ($4/9$, $1/9$ ou $1/36$) e com o vetor de velocidade discreta associado.]
        \vspace{2cm}
    \end{minipage}}
    \fonte{O autor (2026), com base em \textcite{Shan2006}.}
\end{figure}

Duas observações a respeito dessa construção antecipam o desenvolvimento posterior.
A primeira é que a quadratura de três pontos, exata até grau cinco no espaço unidimensional, deixa de sê-lo para os momentos de terceira ordem no espaço multidimensional, o que introduz um erro de ordem $\mathcal{O}(u^3)$ no tensor de fluxo de quantidade de movimento.
A segunda é que os pesos obtidos acima pressupõem temperatura constante e igual à unidade em unidades da rede, hipótese embutida na normalização da gaussiana de referência; quando a temperatura passa a variar de ponto a ponto, como ocorre em escoamentos compressíveis, a mesma construção conduz a pesos dependentes da temperatura, conforme desenvolvido na Seção~\ref{sec:lbm-dupla-distribuicao}.

A função contínua de distribuição $f$ é convertida, correspondentemente, em um conjunto de funções de distribuição direcionais $f_i(\mathbf{x}, t) = w_i f(\mathbf{x}, \mathbf{c}_i, t) / \varpi(\mathbf{c}_i)$, nas quais o peso de quadratura é absorvido de modo que os momentos macroscópicos se reduzam a simples somatórios sobre as direções.
Essas populações são regidas pela equação de Boltzmann discreta (\textit{Discrete Boltzmann Equation}, DBE):
\begin{equation}\label{eq:boltzmann-discreta}
    \frac{\partial f_i}{\partial t} + \mathbf{c}_i \cdot \nabla f_i = -\frac{1}{\tau}(f_i - f_i^{eq})
\end{equation}

O passo seguinte consiste em integrar a \autoref{eq:boltzmann-discreta} ao longo de suas linhas características no espaço e no tempo.
Para um intervalo de tempo $\Delta t$, e considerando que a posição avança de $\mathbf{c}_i \Delta t$, a integração temporal exata assume a forma:
\begin{equation}
    f_i(\mathbf{x} + \mathbf{c}_i \Delta t, t + \Delta t) - f_i(\mathbf{x}, t) = -\int_{0}^{\Delta t} \frac{1}{\tau} \left[ f_i(\mathbf{x} + \mathbf{c}_i s, t + s) - f_i^{eq}(\mathbf{x} + \mathbf{c}_i s, t + s) \right] ds
\end{equation}
A avaliação explícita do lado direito dessa integral demandaria aproximações implícitas e computacionalmente custosas.
Admitindo que a malha espaço-temporal seja suficientemente refinada, o que corresponde a número de Courant unitário, e que $\Delta t$ seja pequeno, a integral pode ser aproximada pelo método de Euler explícito de primeira ordem, no qual o integrando é avaliado no instante inicial $t$.
Essa simplificação conduz à equação de Lattice Boltzmann clássica:
\begin{equation}
    f_i(\mathbf{x} + \mathbf{c}_i \Delta t, t + \Delta t) = f_i(\mathbf{x}, t) - \frac{\Delta t}{\tau} [ f_i(\mathbf{x}, t) - f_i^{eq}(\mathbf{x}, t) ]
\end{equation}
Alternativamente, para elevar a ordem de precisão temporal e preservar a estabilidade, pode-se aplicar a regra do trapézio e, em seguida, introduzir uma mudança de variáveis que define uma função de distribuição modificada, o que também resulta em uma forma explícita do método, com a viscosidade renormalizada por $\Delta t/2$ \parencite{He1997}.
Em todos os casos, o procedimento preserva a subdivisão do algoritmo em duas etapas distintas: a colisão, puramente local, e a propagação (\textit{streaming}), estritamente conservativa.

A dedução precedente estabelece a equação de evolução das populações, mas não fixa o conjunto de velocidades discretas nem a forma explícita da distribuição de equilíbrio adotada no modelo isotérmico.
O modelo bidimensional de nove velocidades, denominado D2Q9, é a formulação mais difundida para a solução de escoamentos incompressíveis isotérmicos, conforme explorado nos trabalhos de \textcite{Qian1992} e \textcite{Chen1998}.
Nessa configuração, as componentes do vetor velocidade discreta $\mathbf{c}_i$ são definidas sobre uma rede quadrada:
\begin{equation}
    \mathbf{c}_i = \begin{cases} 
      (0;\ 0), & i=0 \\ 
      (\pm 1;\ 0)c, \; (0;\ \pm 1)c, & i=1,2,3,4 \\ 
      (\pm 1;\ \pm 1)c, & i=5,6,7,8 
    \end{cases}
\end{equation}
onde a constante $c = \Delta x / \Delta t$ representa a velocidade característica da rede, usualmente adotada como unitária no sistema de unidades da rede.
A correspondência estatística com a função de distribuição contínua é assegurada por pesos de quadratura isotrópicos $w_i$, derivados de modo a anular termos espúrios no tensor de tensões macroscópico e que assumem os valores $w_0 = 4/9$ para a direção de repouso, $w_{1,2,3,4} = 1/9$ para as direções axiais e $w_{5,6,7,8} = 1/36$ para as direções diagonais.

Para simulações tridimensionais, o modelo D3Q19 é frequentemente preferido por equilibrar precisão e custo computacional.
As dezenove velocidades discretas são compostas pelo vetor nulo, seis direções axiais $(\pm 1;\ 0;\ 0)c$, $(0;\ \pm 1;\ 0)c$ e $(0;\ 0;\ \pm 1)c$ e doze direções diagonais nas faces $(\pm 1;\ \pm 1;\ 0)c$, $(\pm 1;\ 0;\ \pm 1)c$ e $(0;\ \pm 1;\ \pm 1)c$.
Os pesos correspondentes são $w_0 = 1/3$ para o repouso, $w_{1..6} = 1/18$ para as direções axiais e $w_{7..18} = 1/36$ para as diagonais.
A \autoref{fig:lbm-redes} apresenta a disposição espacial das velocidades discretas dos modelos D2Q9 e D3Q19.

\begin{figure}[H]
    \centering
    \caption{Conjuntos de velocidades discretas empregados no LBM: (a) modelo bidimensional D2Q9, com nove direções, e (b) modelo tridimensional D3Q19, com dezenove direções.}
    \label{fig:lbm-redes}
    \fbox{\begin{minipage}{0.85\textwidth}
        \centering
        \vspace{2cm}
        [Espaço reservado para figura: dois diagramas lado a lado. Em (a), a rede D2Q9, com o nó central e as nove setas de velocidade discreta numeradas de 0 a 8, indicando as direções de repouso, axiais e diagonais. Em (b), a rede D3Q19, representada por um cubo com o nó central e as dezenove direções, distinguindo as axiais das diagonais das faces.]
        \vspace{2cm}
    \end{minipage}}
    \fonte{Adaptado de \textcite{Chen1998}.}
\end{figure}

A escolha do conjunto de velocidades impacta diretamente a isotropia do método e o custo de memória, conforme comparado na \autoref{tab:lbm-comparacao-redes}.
Redes com maior conectividade, como a D3Q27, oferecem maior estabilidade e precisão em ordens superiores de Hermite, porém ao custo de um aumento significativo no armazenamento e no processamento por nó da malha.

\begin{table}[H]
    \centering
    \caption{Comparação entre conjuntos de velocidades discretas usuais no LBM}
    \label{tab:lbm-comparacao-redes}
    \begin{tabular}{ccccc}
        \toprule
        Modelo & Dimensões & Velocidades & Isotropia (ordem) & Memória por nó (escalares) \\
        \midrule
        D1Q3   & 1         & 3           & 2                 & 3                          \\
        D2Q9   & 2         & 9           & 4                 & 9                          \\
        D3Q19  & 3         & 19          & 4                 & 19                         \\
        D3Q27  & 3         & 27          & 6                 & 27                         \\
        \bottomrule
    \end{tabular}
    \fonte{O autor (2026), com base em \textcite{Qian1992} e \textcite{Chen1998}.}
\end{table}

Para que as simulações recuperem o comportamento incompressível, a distribuição contínua de Maxwell-Boltzmann é truncada por expansão de Taylor, retendo-se os termos até segunda ordem no número de Mach local, $Ma = |\mathbf{u}|/c_s$.
A distribuição discreta de equilíbrio resultante é dada por:
\begin{equation}\label{eq:feq-d2q9}
    f_i^{eq} = w_i \rho \left[ 1 + \frac{\mathbf{c}_i \cdot \mathbf{u}}{c_s^2} + \frac{(\mathbf{c}_i \cdot \mathbf{u})^2}{2c_s^4} - \frac{\mathbf{u} \cdot \mathbf{u}}{2c_s^2} \right]
\end{equation}
onde a velocidade do som na rede, $c_s$, é constante para redes quadradas e assume o valor $c_s = c/\sqrt{3}$ \parencite{Rathnayake2024}.
Os três termos entre colchetes correspondem, respectivamente, aos coeficientes de ordem zero, um e dois da expansão de Hermite da \eqref{eq:expansao-hermite}, de modo que a \eqref{eq:feq-d2q9} é exatamente o truncamento em $\mathcal{N} = 2$ discutido acima.

As grandezas macroscópicas, a massa específica $\rho$ e o vetor velocidade $\mathbf{u}$, são recuperadas por somatórios dos momentos das funções de distribuição em cada nó da malha:
\begin{equation}
    \rho = \sum_{i=0}^8 f_i, \qquad \rho \mathbf{u} = \sum_{i=0}^8 f_i \mathbf{c}_i
\end{equation}
No LBM padrão, a pressão não requer a solução de uma equação de Poisson, como ocorre nas formulações de Navier-Stokes baseadas em métodos de projeção, sendo avaliada localmente por uma equação de estado (\textit{Equation of State}, EOS) linear:
\begin{equation}
    p = c_s^2 \rho
\end{equation}
Essa formulação evidencia que os resolvedores baseados no LBM representam um meio fracamente compressível em condições quase estáticas ou isotérmicas.
O erro induzido por essa compressibilidade artificial é proporcional a $\mathcal{O}(Ma^2)$.
Consequentemente, os resultados hidrodinâmicos do modelo aproximam-se do limite incompressível na medida em que o escoamento apresente números de Mach reduzidos ($Ma \ll 0{,}3$), o que estabelece as fronteiras clássicas de aplicação do método em sua forma básica \parencite{Chen1998}.

Reunidos o conjunto de velocidades discretas, a distribuição de equilíbrio e a equação de evolução, a estrutura computacional do método fica determinada, e é dela que decorre a aptidão do LBM à paralelização massiva.
O ciclo básico compreende quatro etapas: inicialização, colisão, propagação e aplicação das condições de contorno.
O ALGORITMO~\ref{alg:lbm-ciclo} descreve o fluxo lógico de um passo de tempo, cujo comportamento sobre a rede é ilustrado na \autoref{fig:lbm-colisao-propagacao}.

\begin{algorithm}[H]
\caption{Ciclo de execução de um passo de tempo no LBM}\label{alg:lbm-ciclo}
\Entrada{Distribuições $f_i(\mathbf{x}, t)$}
\Saida{Distribuições atualizadas $f_i(\mathbf{x}, t + \Delta t)$}
\Para{cada nó $\mathbf{x}$ da malha}{
    Calcular momentos macroscópicos: $\rho = \sum f_i$ e $\mathbf{u} = \frac{1}{\rho} \sum f_i \mathbf{c}_i$\;
    Calcular a distribuição de equilíbrio $f_i^{eq}(\rho, \mathbf{u})$\;
    \tcp{Colisão}
    $f_i^{coll}(\mathbf{x}, t) = f_i(\mathbf{x}, t) - \frac{1}{\tau} (f_i(\mathbf{x}, t) - f_i^{eq}(\mathbf{x}, t))$\;
}
\Para{cada nó $\mathbf{x}$ da malha}{
    \tcp{Propagação}
    $f_i(\mathbf{x} + \mathbf{c}_i \Delta t, t + \Delta t) = f_i^{coll}(\mathbf{x}, t)$\;
}
Aplicar condições de contorno\;
\Retorne{ $f_i(\mathbf{x}, t + \Delta t)$ }
\end{algorithm}
\fonte{Adaptado de \textcite{Chen1998}.}

\begin{figure}[H]
    \centering
    \caption{Etapas de colisão local e de propagação para os nós vizinhos em uma rede bidimensional D2Q9.}
    \label{fig:lbm-colisao-propagacao}
    \fbox{\begin{minipage}{0.85\textwidth}
        \centering
        \vspace{2cm}
        [Espaço reservado para figura: ilustração das etapas de colisão e de propagação em uma rede D2Q9. No primeiro quadro, as distribuições em um nó antes da colisão; no segundo, as distribuições redistribuídas localmente após a colisão; no terceiro, o deslocamento das distribuições para os nós vizinhos segundo suas direções discretas.]
        \vspace{2cm}
    \end{minipage}}
    \fonte{Adaptado de \textcite{Chen1998}.}
\end{figure}

A separação em dois percursos da malha decorre de uma dependência de dados e não de uma imposição do método, pois a propagação de um nó requer as populações pós-colisão de seus vizinhos.
A colisão permanece inteiramente local e a propagação alcança apenas os vizinhos imediatos, de modo que nenhuma etapa exige reduções globais ou a solução de sistemas lineares acoplados sobre o domínio, o que distingue o método dos resolvedores baseados em projeção de pressão e sustenta a comparação entre tecnologias de paralelização conduzida nos capítulos seguintes.

%---------------------------------------------------------------------------------------

\section{Análise Multiescala de Chapman-Enskog}\label{sec:lbm-chapman-enskog}

A transição formal da escala cinética para a escala hidrodinâmica é estabelecida pela análise de Chapman-Enskog, desenvolvida originalmente por \textcite{Chapman1916} e por \textcite{Enskog1917} para o cálculo dos coeficientes de transporte de gases rarefeitos, sistematizada em \textcite{Chapman1970}.
Essa técnica permite derivar as equações de conservação de massa e de quantidade de movimento a partir da equação de Lattice Boltzmann, tratando as variáveis macroscópicas como variações lentas em relação aos processos de colisão microscópicos.
A análise fundamenta-se na expansão da função de distribuição e das derivadas espaço-temporais em potências de um parâmetro de escala $\epsilon$, proporcional ao número de Knudsen \parencite{Chen1998, Rathnayake2024}.

O ponto de partida é a expansão perturbativa da função de distribuição em torno do estado de equilíbrio local $f_i^{(0)} = f_i^{eq}$:
\begin{equation}
    f_i = f_i^{(0)} + \epsilon f_i^{(1)} + \epsilon^2 f_i^{(2)} + \mathcal{O}(\epsilon^3)
\end{equation}
Para garantir que as variáveis macroscópicas sejam totalmente determinadas pela distribuição de equilíbrio, impõem-se as condições de solubilidade, que estabelecem que as correções de não equilíbrio não contribuem para os momentos primários:
\begin{equation}
    \sum_i f_i^{(k)} = 0, \quad \sum_i f_i^{(k)} \mathbf{c}_i = 0, \quad \sum_i f_i^{(k)} \mathbf{c}_i^2 = 0, \quad k \ge 1
\end{equation}

A análise multiescala introduz diferentes escalas de tempo e de espaço, associadas à advecção e à dissipação viscosa.
As derivadas parciais são expandidas como:
\begin{equation}
    \frac{\partial}{\partial t} = \epsilon \frac{\partial}{\partial t_1} + \epsilon^2 \frac{\partial}{\partial t_2}, \quad \nabla = \epsilon \nabla_1
\end{equation}
Expandindo a equação de Lattice Boltzmann em série de Taylor até segunda ordem e substituindo as expansões de $f_i$ e das derivadas, agrupam-se os termos em potências sucessivas de $\epsilon$.

Na ordem $\epsilon^0$, a equação é identicamente satisfeita pelo estado de equilíbrio.
Na ordem $\epsilon^1$, obtém-se:
\begin{equation}
    \left( \frac{\partial}{\partial t_1} + \mathbf{c}_i \cdot \nabla_1 \right) f_i^{(0)} = -\frac{1}{\tau \Delta t} f_i^{(1)}
\end{equation}
Tomando-se os momentos de ordem zero e um dessa equação, e utilizando as propriedades de conservação, recuperam-se as equações de Euler para fluidos invíscidos:
\begin{equation}
    \frac{\partial \rho}{\partial t_1} + \nabla_1 \cdot (\rho \mathbf{u}) = 0
\end{equation}
\begin{equation}
    \frac{\partial (\rho \mathbf{u})}{\partial t_1} + \nabla_1 \cdot \boldsymbol{\Pi}^{(0)} = 0
\end{equation}
onde o tensor de fluxo de quantidade de movimento no equilíbrio é $\boldsymbol{\Pi}^{(0)} = \sum_i \mathbf{c}_i \mathbf{c}_i f_i^{(0)} = \rho \mathbf{u}\mathbf{u} + p\mathbf{I}$.
A partir da equação de ordem $\epsilon^1$, a correção de primeira ordem $f_i^{(1)}$ pode ser expressa em função do gradiente da distribuição de equilíbrio:
\begin{equation}\label{eq:fi1-solucao}
    f_i^{(1)} = -\tau \Delta t \left( \frac{\partial f_i^{(0)}}{\partial t_1} + \mathbf{c}_i \cdot \nabla_1 f_i^{(0)} \right)
\end{equation}

Na ordem $\epsilon^2$, a equação resultante incorpora os efeitos de transporte viscoso:
\begin{equation}
    \frac{\partial f_i^{(0)}}{\partial t_2} + \left( \frac{\partial}{\partial t_1} + \mathbf{c}_i \cdot \nabla_1 \right) f_i^{(1)} + \frac{\Delta t}{2} \left( \frac{\partial}{\partial t_1} + \mathbf{c}_i \cdot \nabla_1 \right)^2 f_i^{(0)} = -\frac{1}{\tau \Delta t} f_i^{(2)}
\end{equation}
O passo decisivo consiste em tomar o momento de segunda ordem da \eqref{eq:fi1-solucao}, o que fornece a contribuição de não equilíbrio ao tensor de fluxo de quantidade de movimento:
\begin{equation}\label{eq:pi1}
    \boldsymbol{\Pi}^{(1)} = \sum_i \mathbf{c}_i \mathbf{c}_i f_i^{(1)} = -\tau \Delta t \left( \frac{\partial \boldsymbol{\Pi}^{(0)}}{\partial t_1} + \nabla_1 \cdot \mathbf{Q}^{(0)} \right)
\end{equation}
onde $\mathbf{Q}^{(0)} = \sum_i \mathbf{c}_i \mathbf{c}_i \mathbf{c}_i f_i^{(0)}$ é o momento de terceira ordem do equilíbrio.
Substituindo $\boldsymbol{\Pi}^{(0)} = \rho \mathbf{u}\mathbf{u} + \rho c_s^2 \mathbf{I}$, empregando as equações de Euler para eliminar as derivadas temporais e desprezando os termos de ordem $\mathcal{O}(u^3)$, obtém-se
\begin{equation}\label{eq:pi1-final}
    \boldsymbol{\Pi}^{(1)} = -\rho c_s^2 \tau \Delta t \left[ \nabla_1 \mathbf{u} + (\nabla_1 \mathbf{u})^{\mathsf{T}} \right]
\end{equation}
expressão que reproduz a forma newtoniana do tensor de tensões viscosas.
O tensor de fluxo total, $\boldsymbol{\Pi} = \boldsymbol{\Pi}^{(0)} + \epsilon \boldsymbol{\Pi}^{(1)}$, conduz à equação de quantidade de movimento na forma
\begin{equation}
    \frac{\partial (\rho \mathbf{u})}{\partial t} + \nabla \cdot (\rho \mathbf{u}\mathbf{u}) = -\nabla p + \nabla \cdot \boldsymbol{\sigma}, \qquad \boldsymbol{\sigma} = \rho \nu \left[ \nabla \mathbf{u} + (\nabla \mathbf{u})^{\mathsf{T}} \right]
\end{equation}
revelando que a viscosidade cinemática $\nu$ emerge da defasagem entre o equilíbrio e o não equilíbrio:
\begin{equation}\label{eq:viscosidade-tau}
    \nu = c_s^2 \left( \tau - \frac{1}{2} \right) \Delta t
\end{equation}
O termo $-1/2$ não pertence à física do problema, mas à discretização: ele decorre da integração da colisão ao longo do passo de tempo e sua incorporação é o que eleva a precisão do método à segunda ordem no espaço.
A \eqref{eq:viscosidade-tau} também explicita o limite de estabilidade do esquema, uma vez que viscosidades positivas exigem $\tau > \Delta t / 2$, e valores de $\tau$ próximos desse limite tornam a colisão excessivamente rígida.
Essa derivação demonstra que o LBM resolve as equações de Navier-Stokes no limite de baixo número de Mach.

%---------------------------------------------------------------------------------------

\section{Mapeamento de Unidades e Adimensionalização}\label{sec:lbm-unidades}

Uma característica distintiva do LBM é a execução das simulações em um sistema de unidades próprio, denominado unidades da rede (\textit{lattice units}, LU).
Nesse sistema, a distância entre nós da malha ($\Delta x$), o passo de tempo ($\Delta t$) e a massa específica de referência ($\rho_0$) são normalizados para a unidade.
O mapeamento entre o problema físico e a rede é realizado garantindo-se a igualdade de números adimensionais, principalmente o número de Reynolds ($Re = U L / \nu$) e o número de Mach ($Ma = U / c_s$).

O procedimento de conversão inicia-se pela definição dos parâmetros físicos do problema, a saber, o comprimento característico $L$, a velocidade de referência $U$ e a viscosidade cinemática $\nu$.
Em seguida, selecionam-se os parâmetros da rede, isto é, o número de nós $N$ que representam a escala $L$ e a velocidade característica na rede $u_{lb}$, mantida abaixo de $0{,}1$ nas aplicações incompressíveis, para limitar o erro de compressibilidade, e igualada ao número de Mach multiplicado pela velocidade do som da rede quando, como aqui, a compressibilidade é o próprio objeto da simulação.
Determinam-se então os fatores de conversão de comprimento, $C_L = L/N$, e de velocidade, $C_U = U/u_{lb}$, dos quais decorre o fator de conversão de tempo $C_t = C_L / C_U$.
A viscosidade cinemática na rede resulta de $\nu_{lb} = \nu\, C_t / C_L^2$, expressão que se reduz a $\nu_{lb} = u_{lb} N / Re$, e o tempo de relaxação é finalmente obtido por $\tau = \nu_{lb} / c_s^2 + 0{,}5$.

Considera-se um exemplo numérico dimensional para um tubo de choque de comprimento $L = 1{,}0$ m preenchido com ar a $T_0 = 300$ K, cuja massa específica é $\rho_0 = 1{,}2$ kg\,m$^{-3}$ e cuja viscosidade dinâmica é $\mu = 1{,}8 \times 10^{-5}$ Pa\,s, o que corresponde a uma viscosidade cinemática $\nu = 1{,}5 \times 10^{-5}$ m$^2$\,s$^{-1}$, com $\gamma = 1{,}4$, $R = 287$ J\,kg$^{-1}$\,K$^{-1}$ e $Pr = 0{,}71$.
A velocidade do som nesse estado é $a_0 = \sqrt{\gamma R T_0} \approx 347$ m\,s$^{-1}$, e a velocidade de referência correspondente a $Ma = 0{,}25$ é $U = 0{,}25\, a_0 \approx 87$ m\,s$^{-1}$.
O número de Reynolds do problema físico é, portanto, $Re = U L / \nu \approx 5{,}8 \times 10^{6}$, valor característico de escoamentos compressíveis em escala de laboratório.
Adotando-se uma rede com $N = 10^{4}$ nós longitudinais no tubo e uma velocidade na rede $u_{lb} = 0{,}25$, que preserva o número de Mach, a viscosidade na rede resulta em $\nu_{lb} = u_{lb} N / Re \approx 4{,}3 \times 10^{-4}$.
Para o modelo D2Q9, com $c_s^2 = 1/3$, o tempo de relaxação correspondente é $\tau = 3\, \nu_{lb} + 0{,}5 \approx 0{,}5013$.
Esse valor, ainda que admissível, situa-se próximo do limite inferior de estabilidade $\tau = 0{,}5$ estabelecido pela \eqref{eq:viscosidade-tau}, o que expõe a dificuldade característica das simulações compressíveis em elevado número de Reynolds: reproduzir simultaneamente $Re$ e $Ma$ exige tempos de relaxação quase críticos, cuja viabilidade depende dos mecanismos de estabilização discutidos nas seções seguintes.
Refinar a malha atenua essa restrição, pois $\nu_{lb}$ cresce proporcionalmente a $N$, mas eleva na mesma medida o número de passos de tempo necessários para cobrir o intervalo físico simulado, compromisso que caracteriza o mapeamento de unidades no LBM.
A escolha adequada desses parâmetros é essencial para evitar erros de simulação e assegurar que os resultados possam ser convertidos de volta para unidades físicas com precisão \parencite{Rathnayake2024}.

Cabe distinguir esse exemplo da configuração empregada na validação.
O problema do tubo de choque de Sod é definido em grandezas adimensionais, com $\rho_L = 1{,}0$, $p_L = 1{,}0$ e comprimento unitário, e a viscosidade dinâmica adotada, $\mu = 1{,}0 \times 10^{-2}$, é deliberadamente superior à de um gás real nessas escalas.
A ela corresponde um número de Reynolds de referência da ordem de algumas dezenas, escolha que mantém as frentes de onda suficientemente resolvidas para a comparação com a solução analítica e afasta o tempo de relaxação do limite de estabilidade, sem alterar em nada o custo por nó ou o padrão de acesso à memória, que são o objeto da comparação entre tecnologias de paralelização.

%---------------------------------------------------------------------------------------

\section{Limitações do Modelo Clássico em Regime Compressível}\label{sec:lbm-limitacoes-classico}

A análise da seção anterior é válida sob hipóteses restritivas, e é o rompimento dessas hipóteses que motiva toda a formulação adotada neste trabalho.
Quatro limitações do modelo clássico merecem exame detalhado, pois cada uma delas corresponde a uma decisão específica do resolvedor implementado.

A primeira é a equação de estado.
O modelo isotérmico impõe $p = c_s^2 \rho$, com $c_s$ fixado pela rede, o que equivale a uma velocidade do som constante e independente do estado termodinâmico local.
Em um escoamento compressível, ao contrário, a velocidade do som varia com a temperatura segundo $a = \sqrt{\gamma R T}$, e é justamente essa variação que determina a formação e a velocidade de propagação das ondas de choque.
Um resolvedor que mantenha $c_s$ constante não pode representar a razão entre calores específicos $\gamma$ de forma independente e, portanto, não reproduz corretamente as relações de Rankine-Hugoniot.

A segunda é a invariância de Galileu.
A quadratura de três pontos por eixo garante exatidão apenas até os momentos de segunda ordem no espaço multidimensional, de modo que o momento de terceira ordem contém um erro de ordem $\mathcal{O}(u^3)$.
Esse erro, desprezível quando $Ma \ll 0{,}3$, cresce rapidamente com a velocidade e manifesta-se como uma dependência espúria dos resultados em relação ao referencial adotado, comprometendo a captura de estruturas convectadas em alta velocidade.

A terceira é o número de Prandtl.
Com um único tempo de relaxação, a viscosidade e a difusividade térmica ficam vinculadas ao mesmo parâmetro $\tau$, o que fixa $Pr = 1$.
O valor característico para diversos fluidos, como o ar ($Pr = 0{,}71$) adotado neste trabalho, não pode ser imposto sem que se disponha de um segundo parâmetro de relaxação, independente do primeiro.

A quarta é a positividade das populações.
A distribuição de equilíbrio truncada da \eqref{eq:feq-d2q9} é um polinômio, e polinômios não são positivos por construção.
Para velocidades locais elevadas, o termo $-\mathbf{u} \cdot \mathbf{u} / (2c_s^2)$ pode tornar negativas as populações associadas às direções opostas ao escoamento, situação que não possui sentido físico e que, na prática, provoca a divergência da simulação.
Como as descontinuidades de ondas de choque e de condições inicias como as do tubo de choque geram exatamente essas condições de velocidade e de gradiente, a perda de positividade é o modo de falha mais frequente do LBM clássico em regime compressível.

Essas quatro limitações delimitam o que deve ser modificado.
A dependência da velocidade do som em relação à temperatura e a razão de calores específicos arbitrária exigem que a temperatura seja uma variável independente, resolvida por uma segunda função de distribuição.
A invariância de Galileu em alta velocidade exige que o conjunto de velocidades discretas acompanhe o escoamento, o que motiva as redes deslocadas.
O número de Prandtl arbitrário exige um segundo parâmetro de relaxação.
A positividade exige que o equilíbrio deixe de ser um polinômio truncado e volte a ter forma exponencial.
Antes de examinar os instrumentos que atendem a essas exigências, convém situar as diferentes vias pelas quais o LBM foi estendido ao regime compressível, o que permite justificar a formulação adotada neste trabalho e delimitar o alcance dos resultados obtidos.

%---------------------------------------------------------------------------------------

\section{Adaptações do LBM ao Regime Compressível}\label{sec:lbm-adaptacoes-compressivel}

As propostas encontradas na literatura para superar as limitações enumeradas na seção anterior podem ser agrupadas em quatro famílias, distinguidas pela forma como recuperam os momentos de ordem superior que a rede clássica não suporta.
A \autoref{tab:lbm-comparacao-compressivel} reúne essas famílias segundo os critérios que mais importam a um estudo de escalabilidade, a saber, a memória por nó, o alcance da propagação, a possibilidade de ajustar independentemente o número de Prandtl e a razão de calores específicos, e o custo dominante por atualização.

\begin{table}[H]
    \centering
    \caption{Comparação entre as famílias de adaptação do LBM ao regime compressível}
    \label{tab:lbm-comparacao-compressivel}
    \small
    \begin{tabular}{p{2.8cm}p{1.7cm}p{2.4cm}p{1.3cm}p{3.0cm}}
        \toprule
        Abordagem & Populações por nó & Alcance da propagação & $Pr$ e $\gamma$ livres & Custo dominante \\
        \midrule
        LBM clássico isotérmico & $Q$    & vizinhos imediatos & não  & cópia entre vizinhos \\
        Redes multivelocidade   & $2Q$ a $4Q$ & dois a três espaçamentos & sim & memória e zonas de troca \\
        Esquemas híbridos       & $Q + 1$ & estêncil de diferenças finitas & sim & estêncil não local \\
        Dupla distribuição      & $2Q$   & vizinhos imediatos & sim & iteração local e interpolação \\
        Referenciais adaptativos & $2Q$  & vizinhos imediatos & sim & transformação entre referenciais \\
        \bottomrule
    \end{tabular}
    \fonte{O autor (2026), com base em \textcite{Feng2019} e \textcite{Frapolli2015}.}
\end{table}

A primeira família recorre a redes multivelocidade de alta ordem.
Como a insuficiência do modelo clássico decorre do truncamento da quadratura, uma solução direta consiste em elevar a ordem da quadratura de Gauss-Hermite, o que produz conjuntos como o D2Q17 e o D2Q37, com velocidades discretas de módulo superior à unidade \parencite{Watari2003}.
Com momentos exatos até a quarta ou a quinta ordem, uma única função de distribuição passa a reproduzir a equação de energia e a equação de estado dos gases ideais, e a associação dessas redes ao operador entrópico permite alcançar regimes suprassônicos com razão de calores específicos ajustável \parencite{Frapolli2015}.
O custo, entretanto, é duplo: a memória por nó cresce proporcionalmente ao número de velocidades, e a propagação deixa de envolver apenas os vizinhos imediatos, passando a alcançar nós a dois ou três espaçamentos de distância.
Essa perda de localidade compromete justamente a propriedade que torna o LBM atraente em computação de alto desempenho, pois amplia as zonas de troca na decomposição de domínio e complica o tratamento das fronteiras.

A segunda família adota esquemas híbridos.
A rede de baixa conectividade é preservada para a massa e a quantidade de movimento, e a equação de energia, ou de entropia, é resolvida por diferenças finitas, com regularização dos momentos e um sensor de choque para controlar a dissipação \parencite{Feng2019, Guo2020}.
A validade dessa construção para a interação entre onda de choque e camada limite em tubo de choque foi demonstrada por \textcite{Qiu2020}.
A economia de memória é evidente, uma vez que a energia é representada por um único campo escalar em lugar de um conjunto de populações, mas o preço é a coexistência de dois paradigmas numéricos no mesmo resolvedor.
O estêncil de diferenças finitas introduz operações não locais com padrão de acesso distinto ao da propagação, o que exige núcleos de cálculo específicos e reduz a uniformidade da implementação entre as tecnologias de paralelização.

A terceira família, adotada neste trabalho e desenvolvida em detalhe na Seção~\ref{sec:lbm-dupla-distribuicao}, atribui à energia seu próprio conjunto de populações sobre a mesma rede da massa e da quantidade de movimento, e trata a limitação de invariância de Galileu pelo deslocamento do conjunto de velocidades, e não pelo aumento de sua cardinalidade \parencite{Tran2022}.
A memória por nó dobra em relação ao modelo clássico, o que é consideravelmente menos do que a quadruplicação imposta pelo D2Q37, e todas as operações permanecem locais ou restritas aos vizinhos imediatos, mesmo com a propagação interpolada que o deslocamento exige.
Em contrapartida, a solução iterativa que determina o equilíbrio térmico e as tabelas de interpolação elevam o número de operações e o volume de dados por atualização, e a positividade dos pesos da quadratura restringe a temperatura da rede a um intervalo limitado, o que reduz a faixa de variação térmica admissível em uma única simulação.

A quarta família generaliza a ideia do deslocamento ao permitir que cada nó possua seu próprio referencial e sua própria escala de temperatura, procedimento conhecido como partículas sob demanda \parencite{Dorschner2018}.
O método alcança, em princípio, números de Mach arbitrários, pois a velocidade peculiar é mantida pequena em toda parte, mas exige transformar as populações entre referenciais vizinhos a cada propagação, o que reintroduz interpolações e acopla nós adjacentes por operações não triviais.
A formulação empregada neste trabalho pode ser lida como um caso particular dessa construção, no qual o referencial é único e fixado a priori, escolha que preserva a simplicidade do laço de propagação ao preço de exigir que a velocidade de referência seja representativa de todo o domínio.

A leitura conjunta dos quatro casos reunidos na \autoref{tab:lbm-comparacao-compressivel} esclarece a escolha feita.
Como o objeto deste trabalho é a comparação entre tecnologias de paralelização, e não a busca do modelo físico mais preciso, o critério decisivo é que a formulação preserve a estrutura de cálculo que caracteriza o LBM, isto é, colisão local seguida de propagação para vizinhos imediatos, com o mesmo número de operações em todas as células.
A dupla distribuição com rede deslocada é, entre as alternativas examinadas, a que melhor satisfaz esse critério, pois mantém um único paradigma numérico, dispensa estênceis estendidos e concentra o custo adicional em operações aritméticas locais.
Essa concentração é, aliás, de interesse próprio para o estudo, uma vez que eleva a intensidade aritmética do núcleo de cálculo em relação ao LBM isotérmico e afasta o resolvedor, ainda que parcialmente, da condição de limitação estrita pela largura de banda de memória, o que torna a comparação entre tecnologias mais informativa.

Um limite comum a todas as famílias examinadas merece registro.
A simulação de regimes hipersônicos exige considerar efeitos de alta temperatura, como a excitação de modos vibracionais e a dissociação molecular, que a hipótese de gás caloricamente perfeito não contempla.
Embora o modelo de gás ideal seja frequentemente adotado como primeira aproximação, e seja o adotado neste trabalho, extensões do LBM com equações de estado mais complexas e com termos fonte para reações químicas encontram-se em desenvolvimento ativo, o que aproxima o método mesoscópico da física envolvida na reentrada atmosférica \parencite{Guo2020}.

%---------------------------------------------------------------------------------------

\section{Operadores de Colisão e Estabilidade Numérica}\label{sec:lbm-colisao}

Embora o modelo BGK seja amplamente utilizado em razão de sua simplicidade, apresenta limitações intrínsecas, como o número de Prandtl fixo ($Pr = 1$ na formulação padrão) e a ocorrência de instabilidades em regimes de alta velocidade ou de baixa viscosidade.
A raiz dessas instabilidades está no fato de que um único tempo de relaxação impõe a mesma taxa de amortecimento a todos os momentos da função de distribuição.
Os momentos de ordem superior, que não possuem contrapartida hidrodinâmica e apenas acumulam erro de truncamento, são relaxados com a mesma intensidade que o tensor de tensões, do qual depende a viscosidade física, de modo que reduzir a viscosidade para atingir números de Reynolds elevados implica necessariamente reduzir o amortecimento das componentes espúrias.
Como discutido por \textcite{Chen1998}, para superar essas deficiências foram desenvolvidos operadores de colisão mais robustos, entre eles o de múltiplos tempos de relaxação (\textit{Multiple Relaxation Time}, MRT) e o entrópico (\textit{Entropic Lattice Boltzmann Method}, ELBM).

A escolha do operador de colisão não altera o conjunto de velocidades discretas nem os momentos que ele é capaz de representar, de modo que nem a equação de estado isotérmica nem o erro de invariância de Galileu podem ser corrigidos por essa via: ambos decorrem da quadratura, e não da forma como as populações relaxam.
O que o operador governa é a estabilidade do esquema e o número de graus de liberdade disponíveis para o ajuste dos coeficientes de transporte.
Assim, o MRT ataca diretamente a instabilidade associada aos momentos sem contrapartida hidrodinâmica e, ao dispor de taxas independentes, oferece o segundo parâmetro de relaxação exigido pelo número de Prandtl arbitrário, ainda que somente em modelos que já disponham de uma equação de energia.
O ELBM, por sua vez, ataca a perda de positividade, que é o modo de falha mais frequente em alta velocidade.
As duas limitações restantes exigem modificações na própria representação cinética, tratadas na Seção~\ref{sec:lbm-dupla-distribuicao}.

%---------------------------------------------------------------------------------------

\subsection{Múltiplos Tempos de Relaxação}\label{subsec:lbm-mrt}

No modelo MRT, proposto por \textcite{dHumieres1992}, a colisão é realizada no espaço de momentos, e não no espaço das funções de distribuição.
Uma matriz de transformação $\mathbf{M}$ mapeia as funções de distribuição em um conjunto de momentos, $\mathbf{m} = \mathbf{M}\mathbf{f}$, e a colisão ocorre como uma relaxação linear desses momentos em direção a seus estados de equilíbrio:
\begin{equation}\label{eq:colisao-mrt}
    \mathbf{f}^* = \mathbf{f} - \mathbf{M}^{-1}\mathbf{S}[\mathbf{M}\mathbf{f} - \mathbf{m}^{eq}]
\end{equation}
onde $\mathbf{S} = \operatorname{diag}(s_0, s_1, \ldots, s_{Q-1})$ é uma matriz diagonal que contém as taxas de relaxação individuais de cada momento.
Para a rede D2Q9, o vetor de momentos usualmente adotado é
\begin{equation}\label{eq:momentos-d2q9}
    \mathbf{m} = \left( \rho,\ e,\ \varepsilon,\ j_x,\ q_x,\ j_y,\ q_y,\ p_{xx},\ p_{xy} \right)^{\mathsf{T}}
\end{equation}
no qual $\rho$ é a massa específica, $e$ é a energia, $\varepsilon$ é o quadrado da energia, $\mathbf{j} = \rho \mathbf{u}$ é a quantidade de movimento, $q_x$ e $q_y$ são as componentes do fluxo de energia e $p_{xx}$ e $p_{xy}$ são as componentes independentes do tensor de tensões viscosas.
Os momentos de equilíbrio correspondentes, escritos em unidades da rede com massa específica de referência unitária, são
\begin{equation}\label{eq:momentos-equilibrio-d2q9}
    \begin{aligned}
        e^{eq} &= -2\rho + 3 \left( j_x^2 + j_y^2 \right), &\qquad \varepsilon^{eq} &= \rho - 3 \left( j_x^2 + j_y^2 \right), \\
        q_x^{eq} &= -j_x, &\qquad q_y^{eq} &= -j_y, \\
        p_{xx}^{eq} &= j_x^2 - j_y^2, &\qquad p_{xy}^{eq} &= j_x j_y
    \end{aligned}
\end{equation}
conforme a construção de \textcite{Lallemand2000}.
As taxas associadas aos momentos conservados, isto é, à massa específica e à quantidade de movimento, são irrelevantes, pois as respectivas parcelas de não equilíbrio são nulas por construção.
As taxas associadas ao tensor de tensões, isto é, a $p_{xx}$ e a $p_{xy}$, determinam a viscosidade de cisalhamento por meio de $\nu = c_s^2 (1/s_\nu - 1/2) \Delta t$, expressão análoga à \eqref{eq:viscosidade-tau}, e a taxa associada a $e$ determina a viscosidade volumétrica, que no operador BGK permanece presa a um múltiplo fixo da primeira.
As taxas dos momentos restantes, $\varepsilon$, $q_x$ e $q_y$, denominados momentos fantasma por não possuírem significado hidrodinâmico, permanecem livres e podem ser escolhidas de modo a maximizar a margem de estabilidade, conforme a análise de dispersão e de dissipação conduzida por \textcite{Lallemand2000}.
É precisamente essa liberdade que caracteriza o MRT: os momentos fantasma são amortecidos com taxas próximas da unidade, o que os aniquila em poucos passos de tempo, enquanto o tensor de tensões relaxa lentamente, como exige a viscosidade física reduzida dos escoamentos em elevado número de Reynolds.
Atribuindo-se um mesmo valor $s_i = 1/\tau$ a todas as taxas, a \eqref{eq:colisao-mrt} reduz-se ao operador BGK, o que evidencia que este último é um caso particular do primeiro.

Essa flexibilidade permite ajustar de forma independente a viscosidade de cisalhamento e a viscosidade volumétrica, melhorando a estabilidade e a precisão em escoamentos complexos.
O desacoplamento das escalas de relaxação permite que momentos de ordem superior, associados a tensões não físicas, sejam relaxados mais rapidamente, o que estabiliza o resolvedor em números de Reynolds elevados.
\textcite{Benzi1992} destaca que essa abordagem é fundamental em simulações de turbulência, nas quais a dissipação em pequenas escalas deve ser cuidadosamente controlada.
Uma variante econômica é o operador de dois tempos de relaxação (\textit{Two Relaxation Time}, TRT), no qual se distinguem apenas as partes simétrica e antissimétrica das populações em relação à inversão das direções discretas, com taxas $s_+$ e $s_-$.
O produto $\Lambda = (1/s_+ - 1/2)(1/s_- - 1/2)$ controla a posição efetiva das paredes sólidas e a precisão das condições de contorno, o que confere ao TRT boa parte do ganho de estabilidade do MRT com uma fração de seu custo.

O preço do MRT é computacional.
A transformação direta e a inversa da \eqref{eq:colisao-mrt} correspondem, em princípio, a dois produtos matriz-vetor de custo $\mathcal{O}(Q^2)$ por nó e por passo de tempo, além do armazenamento das matrizes $\mathbf{M}$ e $\mathbf{M}^{-1}$.
Implementações eficientes exploram o fato de que as entradas dessas matrizes são inteiras e pequenas, substituindo os produtos por expressões algébricas desenroladas, o que reduz o número de operações mas eleva a pressão sobre os registradores, aspecto crítico em aceleradores de \textit{hardware}, nos quais o número de registradores por fluxo de execução limita diretamente a ocupação dos multiprocessadores.

%---------------------------------------------------------------------------------------

\subsection{Operador Entrópico}\label{subsec:lbm-elbm}

O ELBM parte de uma exigência mais forte do que a simples escolha das taxas de relaxação: que o operador de colisão respeite, já no nível discreto, o teorema H de Boltzmann.
A partir desse teorema e da definição do funcional
\begin{equation}\label{eq:funcional-h-continuo}
    H(t) = \int f \ln f \, \mathrm{d}\boldsymbol{\xi}
\end{equation}
demonstra-se, pela simetria do termo integral de colisão em relação à permuta entre os estados anterior e posterior ao choque binário, que a evolução governada pela equação de Boltzmann satisfaz
\begin{equation}\label{eq:teorema-h}
    \frac{\mathrm{d}H}{\mathrm{d}t} \le 0
\end{equation}
com igualdade se e somente se $f$ coincidir com a distribuição de Maxwell-Boltzmann local \parencite{Chapman1970}.
Como $H$ é, a menos de constante multiplicativa negativa, a entropia por unidade de volume, a \eqref{eq:teorema-h} é a expressão cinética da segunda lei da termodinâmica: as colisões conduzem o gás ao equilíbrio de forma monótona e irreversível, e o equilíbrio é o único estado estacionário admissível.
A relevância numérica do resultado está em que ele fornece uma função de Lyapunov para a dinâmica, isto é, uma grandeza que não pode crescer e cujo mínimo é o equilíbrio, o que exclui a amplificação indefinida de perturbações.
O modelo BGK e a distribuição de equilíbrio truncada não herdam essa propriedade, pois o polinômio da \eqref{eq:feq-d2q9} não é o minimizador de nenhum funcional convexo, e é precisamente essa lacuna que o ELBM se propõe a preencher, ao construir a colisão de modo que um análogo discreto da \eqref{eq:teorema-h} seja satisfeito em cada nó e em cada passo de tempo.

Para tanto, define-se o funcional de entropia discreta
\begin{equation}\label{eq:funcional-h}
    H(\mathbf{f}) = \sum_i f_i \ln \left( \frac{f_i}{w_i} \right)
\end{equation}
cuja monotonicidade ao longo da evolução assegura a estabilidade do esquema \parencite{Karlin1999}.
A distribuição de equilíbrio deixa de ser um polinômio truncado e passa a ser definida como o minimizador de $H$ sujeito às restrições de conservação de massa e de quantidade de movimento, o que garante sua positividade e a consistência termodinâmica do modelo.
Para as redes cujas componentes de velocidade pertencem a $\{-1, 0, +1\}$, esse minimizador admite forma fechada de produto sobre as direções coordenadas \parencite{Ansumali2002}:
\begin{equation}\label{eq:feq-entropico}
    f_i^{eq} = \rho\, w_i \prod_{\alpha=1}^{D} \left( 2 - \sqrt{1 + 3 u_\alpha^2} \right) \left( \frac{2 u_\alpha + \sqrt{1 + 3 u_\alpha^2}}{1 - u_\alpha} \right)^{v_{i\alpha}}
\end{equation}
A expansão dessa expressão em potências da velocidade reproduz, até segunda ordem, o polinômio truncado da \eqref{eq:feq-d2q9}, o que mostra que o equilíbrio clássico é sua aproximação de baixo número de Mach.
A diferença decisiva está fora dessa vizinhança: por ser um produto de fatores estritamente positivos, a \eqref{eq:feq-entropico} nunca produz populações negativas, ao contrário do polinômio, e é essa propriedade que reaparece na \eqref{eq:feq-produto} da formulação compressível adotada neste trabalho, agora com a temperatura local como parâmetro adicional.

A colisão entrópica assume a forma
\begin{equation}\label{eq:colisao-entropica}
    f_i^{col} = f_i + \alpha \beta \left( f_i^{eq} - f_i \right), \qquad \beta = \frac{\Delta t}{2\tau + \Delta t}
\end{equation}
na qual $\beta$ é fixado pela viscosidade desejada e o parâmetro $\alpha$ é determinado, em cada nó e a cada passo de tempo, pela condição de que a colisão não aumente a entropia discreta, ou seja, $H(\mathbf{f} + \alpha (\mathbf{f}^{eq} - \mathbf{f})) = H(\mathbf{f})$.
O valor $\alpha = 2$ corresponde à reflexão exata da distribuição em torno do equilíbrio e recupera a colisão BGK; em regiões de forte não equilíbrio, a raiz da condição entrópica fornece $\alpha < 2$, o que limita o passo da colisão e preserva a positividade das populações.
O mecanismo equivale, portanto, a uma viscosidade adicional que se ativa apenas onde a solução se afasta do equilíbrio, ideia que reaparece, sob outra forma, na relaxação modulada descrita na Seção~\ref{subsec:lbm-sensor}.
\textcite{Frapolli2015} demonstraram que essa construção, associada a redes multivelocidade, estende o LBM a escoamentos compressíveis supersônicos com razão de calores específicos ajustável.
O custo, entretanto, recai sobre a determinação de $\alpha$, obtida por um procedimento de busca de raiz que envolve logaritmos e cujo número de iterações varia de nó a nó, o que introduz divergência entre fluxos de execução e prejudica o desempenho em arquiteturas de paralelismo massivo.

%---------------------------------------------------------------------------------------

\subsection{Comparação entre os Operadores}\label{subsec:lbm-comparacao-operadores}

A \autoref{tab:lbm-operadores} sintetiza os operadores discutidos segundo o número de parâmetros livres, o mecanismo pelo qual obtêm estabilidade e o custo adicional que impõem a cada nó da malha.
Os dois últimos critérios são os que interessam de perto a este trabalho, pois determinam não apenas o número de operações por atualização, mas também a uniformidade do fluxo de execução, propriedade decisiva para o desempenho em aceleradores.

\begin{table}[H]
    \centering
    \caption{Comparação entre operadores de colisão quanto a parâmetros livres, mecanismo de estabilização e custo adicional por nó}
    \label{tab:lbm-operadores}
    \begin{tabular}{lcp{4.0cm}p{3.6cm}}
        \toprule
        Operador & Parâmetros livres & Mecanismo de estabilização & Custo adicional por nó \\
        \midrule
        BGK              & $1$          & nenhum                                          & nulo                                   \\
        TRT              & $2$          & separação entre partes simétrica e antissimétrica & poucas operações adicionais            \\
        MRT              & até $Q$      & amortecimento dos momentos fantasma             & dois produtos matriz-vetor             \\
        ELBM             & $1$ e $\alpha$ & teorema H discreto                            & busca de raiz com logaritmos           \\
        Regularização    & $1$          & descarte dos momentos não suportados            & projeções em polinômios de Hermite     \\
        Presente trabalho & $2$         & equilíbrio entrópico e relaxação modulada       & iteração de Newton-Raphson local       \\
        \bottomrule
    \end{tabular}
    \fonte{O autor (2026), com base em \textcite{dHumieres1992} e \textcite{Karlin1999}.}
\end{table}

O miniaplicativo desenvolvido não emprega o MRT nem a busca do parâmetro entrópico.
Adota-se um único fator de relaxação para cada uma das duas distribuições, modulado localmente pelo indicador de não equilíbrio da Seção~\ref{subsec:lbm-sensor}, e o equilíbrio da energia é construído por maximização de entropia, conforme a Seção~\ref{subsec:lbm-equilibrio-g}, o que transfere ao resolvedor a garantia de positividade característica do ELBM sem exigir a solução da condição entrópica em cada nó.
A escolha é deliberada e decorre do objeto de estudo: como o interesse recai sobre a comparação entre tecnologias de paralelização, privilegiam-se operadores cujo número de operações por nó seja previsível e cujo fluxo de execução seja idêntico em todas as células, de modo que as diferenças de desempenho observadas no Capítulo~\ref{ch:resultados} possam ser atribuídas ao modelo de execução e não a desequilíbrios de carga introduzidos pelo próprio operador de colisão.

%---------------------------------------------------------------------------------------

\section{Regularização de Momentos}\label{sec:lbm-regularizacao}

Em regimes de escoamento compressível, a estabilidade do LBM é frequentemente comprometida por erros de truncamento na quadratura de Hermite.
A origem do problema está na assimetria entre o que a rede representa no equilíbrio e o que ela carrega fora dele.
A distribuição de equilíbrio da \eqref{eq:feq-d2q9} pertence, por construção, ao espaço gerado pelos polinômios de Hermite de ordem até dois, ao passo que a parcela de não equilíbrio $f_i^{neq} = f_i - f_i^{eq}$, produzida a cada passo de tempo pela propagação, ocupa todo o espaço de dimensão $Q$ das populações.
Das nove componentes de $f_i^{neq}$ na rede D2Q9, apenas as três do tensor $\boldsymbol{\Pi}^{(1)}$ possuem significado hidrodinâmico; as demais são os momentos fantasma já mencionados, que a colisão BGK relaxa com a mesma taxa lenta atribuída à viscosidade e que, por isso, se acumulam em regiões de gradiente acentuado.
A \autoref{fig:lbm-regularizacao} representa essa decomposição do espaço de momentos.

\begin{figure}[H]
    \centering
    \caption{Decomposição do espaço de momentos da rede D2Q9 e efeito da regularização sobre a parcela de não equilíbrio.}
    \label{fig:lbm-regularizacao}
    \fbox{\begin{minipage}{0.85\textwidth}
        \centering
        \vspace{2cm}
        [Espaço reservado para figura: dois quadros lado a lado representando o espaço de momentos da rede D2Q9. No quadro da esquerda, o vetor de não equilíbrio decomposto em duas parcelas, uma no subespaço hidrodinâmico, associado ao tensor de tensões, e outra no subespaço dos momentos fantasma. No quadro da direita, o mesmo vetor após a regularização, com a parcela fantasma anulada e a projeção hidrodinâmica preservada, indicando que apenas o conteúdo suportado pela rede é reconstruído.]
        \vspace{2cm}
    \end{minipage}}
    \fonte{O autor (2026), com base em \textcite{Latt2006}.}
\end{figure}

A técnica de regularização consiste em descartar essas componentes antes da colisão, substituindo $f_i^{neq}$ por sua projeção no espaço de polinômios de Hermite efetivamente suportado pela rede \parencite{Latt2006}.
O procedimento inicia-se pela medição dos coeficientes de não equilíbrio,
\begin{equation}\label{eq:coeficientes-neq}
    \mathbf{a}_{neq}^{(n)} = \sum_i f_i^{neq}\, \boldsymbol{\mathcal{H}}^{(n)}(\mathbf{c}_i)
\end{equation}
e prossegue com a reconstrução da parcela de não equilíbrio a partir apenas dos coeficientes retidos:
\begin{equation}\label{eq:regularizacao}
    f_i^{neq, reg} = w_i \sum_{n=2}^{\mathcal{N}} \frac{1}{n!\, c_s^{2n}}\, \mathbf{a}_{neq}^{(n)} : \boldsymbol{\mathcal{H}}^{(n)}(\mathbf{c}_i)
\end{equation}
Retendo-se apenas o termo de segunda ordem, para o qual $\mathbf{a}_{neq}^{(2)} = \boldsymbol{\Pi}^{(1)}$, a \eqref{eq:regularizacao} reduz-se a
\begin{equation}\label{eq:regularizacao-segunda-ordem}
    f_i^{neq, reg} = \frac{w_i}{2 c_s^4}\, \Pi_{\alpha\beta}^{(1)} \left( c_{i\alpha} c_{i\beta} - c_s^2 \delta_{\alpha\beta} \right)
\end{equation}
e a colisão passa a ser escrita na forma $f_i^{col} = f_i^{eq} + (1 - 1/\tau)\, f_i^{neq, reg}$.
O efeito é imediato: as componentes fantasma são anuladas a cada passo de tempo, em vez de amortecidas lentamente, e a estabilidade deixa de depender do valor de $\tau$ na mesma medida.

A regularização recursiva (\textit{Recursive Regularization}, RR) estende a soma da \eqref{eq:regularizacao} a ordens superiores sem precisar medi-las nas populações, o que seria impossível com exatidão em uma rede de baixa conectividade.
Em lugar disso, os coeficientes de ordem mais alta são obtidos recursivamente dos de ordem inferior e da velocidade macroscópica \parencite{Coreixas2017}:
\begin{equation}\label{eq:regularizacao-recursiva}
    a_{neq, \alpha\beta\gamma}^{(3)} = u_\alpha a_{neq, \beta\gamma}^{(2)} + u_\beta a_{neq, \alpha\gamma}^{(2)} + u_\gamma a_{neq, \alpha\beta}^{(2)}
\end{equation}
relação que decorre da própria estrutura da expansão de Chapman-Enskog e que se prolonga, de forma análoga, às ordens seguintes.
Como todos os coeficientes passam a derivar de $\boldsymbol{\Pi}^{(1)}$, o esquema recupera parte dos termos de ordem $\mathcal{O}(u^3)$ perdidos no truncamento e atenua, por essa via, o erro de invariância de Galileu discutido na Seção~\ref{sec:lbm-limitacoes-classico}, ainda que sem eliminá-lo, uma vez que a quadratura permanece a mesma \parencite{Guo2020}.

A versão híbrida (\textit{Hybrid Recursive Regularization}, HRR) reconhece que a medição de $\boldsymbol{\Pi}^{(1)}$ nas próprias populações é pouco confiável junto às descontinuidades e a substitui por uma combinação com a taxa de deformação avaliada por diferenças finitas:
\begin{equation}\label{eq:hrr}
    \mathbf{a}_{neq}^{(2)} = \varsigma\, \boldsymbol{\Pi}^{(1)} - (1 - \varsigma)\, \rho\, c_s^2 \tau \Delta t \left[ \nabla \mathbf{u} + (\nabla \mathbf{u})^{\mathsf{T}} \right], \qquad 0 \le \varsigma \le 1
\end{equation}
onde o parâmetro $\varsigma$ pondera as duas estimativas e controla, na prática, a dissipação adicional introduzida.
Essa técnica, validada por \textcite{Qiu2020}, permite o uso de redes de baixa conectividade, como D2Q9 ou D3Q19, em números de Mach elevados, para os quais tradicionalmente seriam exigidas redes de ordem superior.
O preço é o mesmo dos esquemas híbridos discutidos na Seção~\ref{sec:lbm-adaptacoes-compressivel}, ou seja, a introdução de um estêncil de diferenças finitas cujo padrão de acesso difere do da propagação.

A regularização RR e a HRR relacionam-se ao MRT no sentido de que ambas operam sobre os momentos do escoamento, porém a regularização impõe uma estrutura polinomial estrita que reforça a consistência com a expansão de Hermite original.
Enquanto o MRT atribui taxas de relaxação distintas a momentos que continuam presentes na distribuição, a regularização simplesmente descarta os momentos não suportados pela rede, o que a torna mais econômica em memória, por dispensar as matrizes de transformação, e mais previsível em termos de fluxo de execução, propriedade relevante para as implementações voltadas a aceleradores de \textit{hardware}.

Em escoamentos hipersônicos, a severidade das descontinuidades torna necessária estabilização adicional.
Além das técnicas de regularização, empregam-se esquemas como o de Lax-Wendroff aplicado ao LBM ou interpolações a montante (\textit{upwind}), entre elas o WENO, na etapa de propagação.
Essas técnicas introduzem uma dissipação numérica controlada que inibe as oscilações de Gibbs nas proximidades dos choques, permitindo a simulação estável de números de Mach elevados ($Ma > 5$) \parencite{Tran2022}.

O miniaplicativo desenvolvido não emprega a projeção de Hermite dessas formulações.
A estabilidade é obtida por dois mecanismos distintos, ambos propostos por \textcite{Tran2022} e detalhados na Seção~\ref{sec:lbm-dupla-distribuicao}: a construção iterativa de um equilíbrio térmico de forma exponencial, positivo por construção, e uma relaxação cujo valor local depende do afastamento medido em relação ao equilíbrio.
O efeito pretendido é o mesmo das técnicas de regularização, isto é, conter as oscilações de Gibbs nas proximidades das descontinuidades, mas o custo por nó concentra-se em uma solução iterativa local, e não em produtos matriciais, o que altera de maneira significativa o perfil de acesso à memória discutido no Capítulo~\ref{ch:implementacao-computacional}.

%---------------------------------------------------------------------------------------

\section{Dupla Distribuição e Redes Deslocadas}\label{sec:lbm-dupla-distribuicao}

A formulação apresentada nesta seção é a que se encontra implementada no miniaplicativo e responde, uma a uma, às quatro limitações enumeradas na Seção~\ref{sec:lbm-limitacoes-classico}.
Ela foi proposta por \textcite{Tran2022} para escoamentos compressíveis em elevado número de Reynolds e posteriormente detalhada e estendida por \textcite{Rathnayake2024} para o regime supersônico.
Seus quatro elementos constitutivos são a dupla função de distribuição (\textit{Double Distribution Function}, DDF), as redes deslocadas (\textit{shifted lattices}), o equilíbrio térmico obtido por maximização de entropia e a relaxação dependente do não equilíbrio local.
A correspondência entre esses elementos e as limitações que os motivam é direta e representada na \autoref{tab:lbm-limitacoes-solucoes},

\begin{table}[H]
    \centering
    \caption{Correspondência entre as limitações do modelo clássico e os elementos da formulação adotada}
    \label{tab:lbm-limitacoes-solucoes}
    \begin{tabular}{p{4.2cm}p{4.6cm}p{3.2cm}}
        \toprule
        Limitação do modelo clássico & Elemento que a resolve \\
        \midrule
        Equação de estado isotérmica & Segunda distribuição para a energia \\
        Erro de invariância de Galileu & Rede deslocada pela velocidade de referência \\
        Número de Prandtl fixo & Segundo fator de relaxação com quase equilíbrio \\
        Perda de positividade & Equilíbrios de forma exponencial e fatorada \\
        \bottomrule
    \end{tabular}
    \fonte{O autor (2026), com base em \textcite{Tran2022}.}
\end{table}

%---------------------------------------------------------------------------------------

\subsection{Equações de Evolução das Duas Distribuições}\label{subsec:lbm-ddf-evolucao}

A primeira decisão consiste em desacoplar a energia da massa e da quantidade de movimento, atribuindo-lhes conjuntos distintos de populações.
O conjunto $f_i$ transporta a massa específica e a quantidade de movimento, ao passo que o conjunto $g_i$ transporta a energia total.
Suas equações de evolução, na forma proposta por \textcite{Tran2022}, são
\begin{equation}\label{eq:ddf-f}
    f_i(\mathbf{x} + \mathbf{c}_i \Delta t, t + \Delta t) = f_i(\mathbf{x}, t) + \omega \left[ f_i^{eq}(\mathbf{x}, t) - f_i(\mathbf{x}, t) \right]
\end{equation}
\begin{equation}\label{eq:ddf-g}
    g_i(\mathbf{x} + \mathbf{c}_i \Delta t, t + \Delta t) = g_i(\mathbf{x}, t) + \omega \left[ g_i^{eq}(\mathbf{x}, t) - g_i(\mathbf{x}, t) \right] + (\omega - \omega_T) \left[ g_i^{*}(\mathbf{x}, t) - g_i^{eq}(\mathbf{x}, t) \right]
\end{equation}
onde $\omega = 1/\tau$ e $\omega_T = 1/\tau_T$ são os fatores de relaxação associados à quantidade de movimento e à energia, e $g_i^{*}$ é uma distribuição de quase equilíbrio detalhada na Seção~\ref{subsec:lbm-transporte}.
A \eqref{eq:ddf-f} preserva a estrutura da equação de Lattice Boltzmann clássica.
A \eqref{eq:ddf-g} distingue-se pelo termo adicional, cuja presença é o que permite atribuir a $\omega_T$ um valor independente de $\omega$ sem violar a conservação da energia total, resolvendo assim a limitação do número de Prandtl fixo.

As grandezas macroscópicas são recuperadas pelos momentos
\begin{equation}\label{eq:ddf-momentos}
    \rho = \sum_i f_i, \qquad \rho \mathbf{u} = \sum_i \mathbf{c}_i f_i, \qquad 2 \rho E = \sum_i g_i
\end{equation}
com a energia total por unidade de massa dada por $E = c_v T + |\mathbf{u}|^2/2$, onde $c_v = R/(\gamma - 1)$ é o calor específico a volume constante.
A temperatura local resulta, portanto, da inversão direta dessa relação,
\begin{equation}\label{eq:temperatura-local}
    T = \frac{1}{c_v} \left( \frac{1}{2\rho} \sum_i g_i - \frac{|\mathbf{u}|^2}{2} \right)
\end{equation}
e a pressão passa a ser calculada pela equação de estado dos gases ideais, $p = \rho R T$, e não mais pela relação isotérmica $p = c_s^2 \rho$.
Com isso, a velocidade do som local torna-se $a = \sqrt{\gamma R T}$ e passa a variar de nó a nó, o que é a condição necessária para representar corretamente a propagação das ondas de choque.

A razão de calores específicos merece comentário à parte, por ser o parâmetro que o modelo clássico não consegue ajustar.
Na formulação de dupla distribuição, $\gamma$ não é imposta pela rede, e sim pelo calor específico com que a temperatura é convertida em energia interna, uma vez que $c_v = R/(\gamma - 1)$ equivale a atribuir ao gás um número de graus de liberdade igual a $2 c_v / R$.
A rede fornece apenas os graus de liberdade translacionais associados às $D$ direções do espaço, ao passo que os demais, representativos dos modos internos da molécula, são contabilizados na própria definição de $E$ e não exigem populações adicionais.
Por esse motivo, o mesmo conjunto de nove velocidades atende tanto ao gás monoatômico quanto ao ar, para o qual se adota $\gamma = 1{,}4$ neste trabalho.
O número de Mach local, empregado adiante na análise dos resultados, resulta então de $Ma = |\mathbf{u}| / a$, com $a$ avaliado no estado termodinâmico de cada nó, e não mais em relação à velocidade do som fixa da rede.

%---------------------------------------------------------------------------------------

\subsection{Redes Deslocadas e Propagação Interpolada}\label{subsec:lbm-redes-deslocadas}

Em redes deslocadas, o conjunto de velocidades discretas é transladado por uma velocidade de referência constante $\mathbf{U}_{ref}$:
\begin{equation}\label{eq:rede-deslocada}
    \mathbf{c}_i = \mathbf{v}_i + \mathbf{U}_{ref}
\end{equation}
onde $\mathbf{v}_i$ designa as velocidades da rede padrão, de componentes em $\{-1, 0, +1\}$.
O efeito dessa translação é substituir a velocidade macroscópica $\mathbf{u}$ pela velocidade particular $\tilde{\mathbf{u}} = \mathbf{u} - \mathbf{U}_{ref}$ em todas as expressões de equilíbrio.
Como os erros de truncamento e a perda de positividade escalam com potências dessa velocidade, escolher $\mathbf{U}_{ref}$ próxima da velocidade média do escoamento reduz simultaneamente o erro de invariância de Galileu e o risco de populações negativas.
Em termos geométricos, o deslocamento centra a quadratura na região do espaço de velocidades em que a função de distribuição efetivamente se concentra, conforme ilustrado na \autoref{fig:lbm-rede-deslocada}.
No miniaplicativo, $\mathbf{U}_{ref}$ é um parâmetro de entrada, fixado em $U_x = 0{,}25$ para o tubo de choque de Sod, valor correspondente à velocidade média do escoamento no intervalo simulado.

\begin{figure}[H]
    \centering
    \caption{Comparação entre uma rede padrão centrada na origem e uma rede deslocada pela velocidade de referência $\mathbf{U}_{ref}$, que otimiza a representação da função de distribuição em escoamentos de alta velocidade.}
    \label{fig:lbm-rede-deslocada}
    \fbox{\begin{minipage}{0.85\textwidth}
        \centering
        \vspace{2cm}
        [Espaço reservado para figura: comparação, no espaço de velocidades, entre uma rede padrão centrada na origem e uma rede deslocada pela velocidade de referência. Em cada caso, deve ser sobreposta a função de distribuição do escoamento supersônico, evidenciando que a rede deslocada mantém a distribuição no interior do conjunto de velocidades discretas.]
        \vspace{2cm}
    \end{minipage}}
    \fonte{Adaptado de \textcite{Tran2022}.}
\end{figure}

O ganho obtido com o deslocamento pode ser estimado, já que o erro de invariância de Galileu de uma rede de três pontos por eixo passa a ser de ordem $\mathcal{O}(\tilde{u}^3)$, e não mais $\mathcal{O}(u^3)$, uma vez que é a velocidade relativa, e não a velocidade absoluta, que comparece nas expressões de equilíbrio.
Um escoamento uniformemente convectivo com velocidade igual à de referência é, portanto, representado exatamente, e o erro remanescente mede apenas o afastamento local em relação ao referencial adotado.
Daí decorre o critério de escolha: o ganho é tanto maior quanto mais estreita for a faixa de velocidades presentes no domínio, e a melhor escolha é a que centra o referencial nessa faixa.
Aí reside também a limitação do procedimento, pois um único referencial não serve igualmente bem a regiões com velocidades muito distintas, situação que motivou os referenciais adaptativos mencionados na Seção~\ref{sec:lbm-adaptacoes-compressivel}.
No tubo de choque, a faixa de velocidades é estreita o bastante para que um referencial único seja adequado, o que faz do problema um caso favorável à formulação adotada.

O deslocamento traz consigo uma consequência algorítmica de peso.
Como $\mathbf{U}_{ref}$ não é, em geral, um número inteiro, o ponto de partida $\mathbf{x} - \mathbf{c}_i \Delta t$ deixa de coincidir com um nó da malha, e a propagação, que no LBM clássico é uma simples cópia de valores entre vizinhos, passa a exigir interpolação.
Adota-se a ponderação pelo inverso da distância (\textit{Inverse Distance Weighting}, IDW) sobre os $2 \cdot D$ nós que cercam o ponto de partida:
\begin{equation}\label{eq:propagacao-idw}
    f_i(\mathbf{x}, t + \Delta t) = \sum_{k=1}^{2 \cdot D} \varphi_k\, f_i^{col}(\mathbf{x}_k, t), \qquad \varphi_k = \frac{d_k^{-q}}{\sum_{m=1}^{2 \cdot D} d_m^{-q}}
\end{equation}
onde $\mathbf{x}_k$ são os nós que cercam o ponto de partida, $d_k$ é a distância entre esse ponto e o nó $k$, $q$ é o expoente da ponderação e $f_i^{col}$ designa as populações após a colisão.
Quando o ponto de partida coincide com um nó, atribui-se peso unitário a esse nó e nulo aos demais, o que recupera a propagação exata do LBM clássico.
A \autoref{fig:lbm-propagacao-idw} representa esse procedimento.
Como os pesos $\varphi_k$ dependem apenas da geometria da malha e de $\mathbf{U}_{ref}$, ambos invariantes ao longo da simulação, eles são calculados uma única vez e armazenados em tabelas, decisão de implementação retomada no Capítulo~\ref{ch:implementacao-computacional}.
Cabe registrar que essa interpolação introduz difusão numérica, contribuição que se soma à dissipação física e que ajuda a explicar a espessura observada das frentes de choque discutida no Capítulo~\ref{ch:resultados}.

\begin{figure}[H]
    \centering
    \caption{Propagação em rede deslocada: o ponto de partida não coincide com um nó da malha, e a população transportada resulta da ponderação pelo inverso da distância dos quatro nós vizinhos.}
    \label{fig:lbm-propagacao-idw}
    \fbox{\begin{minipage}{0.85\textwidth}
        \centering
        \vspace{2cm}
        [Espaço reservado para figura: detalhe de uma malha cartesiana com quatro nós vizinhos formando uma célula. No interior da célula, um ponto assinalado como origem da propagação, situado em posição não coincidente com os nós por efeito do deslocamento da rede, ligado por segmentos tracejados aos quatro nós, com as distâncias e os pesos correspondentes indicados. Uma seta deve partir desse ponto em direção ao nó de destino, indicando o sentido da propagação.]
        \vspace{2cm}
    \end{minipage}}
    \fonte{Adaptado de \textcite{Rathnayake2024}.}
\end{figure}

%---------------------------------------------------------------------------------------

\subsection{Pesos Dependentes da Temperatura e Equilíbrio de Massa}\label{subsec:lbm-equilibrio-f}

Retomando a construção da Seção~\ref{sec:lbm-deducao}, mas admitindo agora que a temperatura varie de ponto a ponto, os pesos da quadratura deixam de ser constantes e passam a acompanhar o estado térmico local.
Para uma rede cujas componentes pertencem a $\{-1, 0, +1\}$, os pesos assumem a forma de produto sobre as direções coordenadas:
\begin{equation}\label{eq:pesos-temperatura}
    W_i(T) = \prod_{\alpha=1}^{D} w(v_{i\alpha}, T), \qquad w(0, T) = 1 - T, \qquad w(\pm 1, T) = \frac{T}{2}
\end{equation}
Esses pesos satisfazem, por construção, as condições de isotropia $\sum_i W_i = 1$, $\sum_i W_i v_{i\alpha} = 0$ e $\sum_i W_i v_{i\alpha} v_{i\beta} = T \delta_{\alpha\beta}$, sendo $\delta_{\alpha\beta}$ o delta de Kronecker.
A última dessas relações revela o significado da temperatura em unidades da rede: ela é a variância do conjunto de velocidades discretas.
Substituindo $T = 1/3$ na \eqref{eq:pesos-temperatura}, obtém-se $w(0) = 2/3$ e $w(\pm 1) = 1/6$, cujos produtos reproduzem exatamente os pesos $4/9$, $1/9$ e $1/36$ do modelo clássico.
O modelo isotérmico é, portanto, o caso particular desta formulação em que a temperatura da rede é congelada no valor $c_s^2$.
A positividade dos pesos restringe a temperatura em unidades da rede ao intervalo $0 < T < 1$, o que estabelece, na prática, a faixa de variação térmica admissível e, por consequência, a escala de temperatura adotada no mapeamento de unidades.

A distribuição de equilíbrio de massa e de quantidade de movimento é construída como um produto de fatores independentes por direção coordenada:
\begin{equation}\label{eq:feq-produto}
    f_i^{eq} = \rho \prod_{\alpha=1}^{D} \Phi(v_{i\alpha}, \tilde{u}_\alpha, T)
\end{equation}
com $\tilde{u}_\alpha = u_\alpha - U_{ref,\alpha}$ e
\begin{equation}\label{eq:feq-fatores}
    \Phi(0) = 1 - \left( \tilde{u}_\alpha^2 + T \right), \qquad
    \Phi(\pm 1) = \frac{\pm \tilde{u}_\alpha + \tilde{u}_\alpha^2 + T}{2}
\end{equation}
A verificação direta dos momentos mostra que essa forma satisfaz $\sum_i f_i^{eq} = \rho$, $\sum_i \mathbf{c}_i f_i^{eq} = \rho \mathbf{u}$ e $\sum_i \mathbf{c}_i \mathbf{c}_i f_i^{eq} = \rho \mathbf{u}\mathbf{u} + \rho T \mathbf{I}$, o que recupera o tensor de pressão correto com a pressão termodinâmica $p = \rho T$ em unidades da rede.
A vantagem da forma fatorada é dupla.
Do ponto de vista numérico, cada fator é exato em sua direção, de modo que o erro remanescente na invariância de Galileu é de ordem $\mathcal{O}(Ma^4)$, contra $\mathcal{O}(Ma^3)$ do polinômio truncado da \eqref{eq:feq-d2q9}.
Do ponto de vista computacional, o custo cresce com $\mathcal{O}(D)$ por direção, e não com o número de termos de uma expansão polinomial, o que se traduz em um núcleo de cálculo compacto e sem ramificações de fluxo.

%---------------------------------------------------------------------------------------

\subsection{Equilíbrio de Energia por Maximização de Entropia}\label{subsec:lbm-equilibrio-g}

A distribuição de equilíbrio da energia não admite uma forma fatorada análoga.
Em vez de postular um polinômio, que reintroduziria o problema da positividade, determina-se $g_i^{eq}$ como a distribuição que extremiza a entropia discreta
\begin{equation}\label{eq:entropia-discreta}
    \mathcal{S} = \sum_i g_i \ln \left( \frac{g_i}{W_i(T)} \right)
\end{equation}
sujeita às restrições de conservação da energia total e do fluxo de energia:
\begin{equation}\label{eq:restricoes-energia}
    \sum_i g_i^{eq} = 2 \rho E, \qquad \sum_i \mathbf{c}_i g_i^{eq} = 2 \rho \mathbf{u} \left( E + T \right)
\end{equation}
A segunda restrição merece atenção, pois o fluxo de energia de um gás ideal compreende não apenas o transporte convectivo da energia total, mas também o trabalho das forças de pressão, e é o termo $\rho \mathbf{u} T$ que o representa.
Impondo as restrições por multiplicadores de Lagrange, a solução do problema variacional assume forma exponencial:
\begin{equation}\label{eq:geq-entropico}
    g_i^{eq} = \rho\, W_i(T) \exp \left( \chi + \boldsymbol{\lambda} \cdot \mathbf{c}_i \right)
\end{equation}
onde $\chi$ é o multiplicador associado à energia total e $\boldsymbol{\lambda}$ o vetor de multiplicadores associados ao fluxo de energia.
Por ser uma exponencial multiplicada por pesos positivos, $g_i^{eq}$ é positiva para quaisquer valores dos multiplicadores, o que elimina por construção o principal modo de falha do LBM clássico em alta velocidade.
O preço dessa garantia é a perda da forma fechada: os multiplicadores não possuem expressão analítica e devem ser determinados numericamente em cada nó e a cada passo de tempo.

Convém introduzir a notação $e_i(\boldsymbol{\lambda}) = W_i(T) \exp(\boldsymbol{\lambda} \cdot \mathbf{c}_i)$ e a soma de normalização $Z(\boldsymbol{\lambda}) = \sum_i e_i$.
A primeira das restrições da \eqref{eq:restricoes-energia} determina $\chi$ de forma explícita, pois basta escalar a distribuição para que sua soma coincida com $2\rho E$.
Resta, assim, um sistema de $D$ equações não lineares nas $D$ componentes de $\boldsymbol{\lambda}$, cujo resíduo é
\begin{equation}\label{eq:residuo-newton}
    r_\alpha(\boldsymbol{\lambda}) = \frac{1}{Z} \sum_i c_{i\alpha} e_i - \frac{u_\alpha (E + T)}{E} = 0
\end{equation}
A matriz jacobiana desse sistema possui interpretação estatística direta:
\begin{equation}\label{eq:jacobiana-newton}
    J_{\alpha\beta} = \frac{\partial r_\alpha}{\partial \lambda_\beta} = \frac{1}{Z} \sum_i c_{i\alpha} c_{i\beta} e_i - \left( \frac{1}{Z} \sum_i c_{i\alpha} e_i \right) \left( \frac{1}{Z} \sum_i c_{i\beta} e_i \right)
\end{equation}
isto é, trata-se da matriz de covariância das velocidades discretas sob a distribuição normalizada $e_i / Z$.
Por essa razão, $\mathbf{J}$ é simétrica e definida positiva sempre que ao menos duas direções tenham peso não nulo, propriedade que assegura a boa condição do sistema e a existência da correção de Newton.
A iteração assume a forma
\begin{equation}\label{eq:atualizacao-newton}
    \boldsymbol{\lambda}^{(k+1)} = \boldsymbol{\lambda}^{(k)} - \alpha_k\, \mathbf{J}^{-1} \mathbf{r}\!\left( \boldsymbol{\lambda}^{(k)} \right)
\end{equation}
onde $\alpha_k$ é um fator de amortecimento que decresce a cada iteração, adotado no miniaplicativo como $\alpha_k = 2^{-k}$, o que evita passos excessivos nas células situadas sobre as descontinuidades.
O ALGORITMO~\ref{alg:lbm-newton} resume o procedimento.

\begin{algorithm}[H]
\caption{Determinação dos multiplicadores de Lagrange do equilíbrio térmico}\label{alg:lbm-newton}
\Entrada{Massa específica $\rho$, velocidade $\mathbf{u}$, temperatura $T$ e multiplicadores $\boldsymbol{\lambda}$ do passo anterior}
\Saida{Distribuição de equilíbrio de energia $g_i^{eq}$}
$\alpha \gets 1$\;
\Para{$k = 1$ até $k_{max}$}{
    Avaliar $e_i = W_i(T) \exp(\boldsymbol{\lambda} \cdot \mathbf{c}_i)$ e $Z = \sum_i e_i$\;
    Montar o resíduo $\mathbf{r}$ pela \eqref{eq:residuo-newton} e a jacobiana $\mathbf{J}$ pela \eqref{eq:jacobiana-newton}\;
    $\boldsymbol{\lambda} \gets \boldsymbol{\lambda} - \alpha\, \mathbf{J}^{-1} \mathbf{r}$\;
    $\alpha \gets \alpha / 2$\;
}
Reavaliar $e_i$ e $Z$ com o valor final de $\boldsymbol{\lambda}$\;
$g_i^{eq} \gets \dfrac{2 \rho E}{Z}\, e_i$\;
\Retorne{$g_i^{eq}$}
\end{algorithm}
\fonte{Adaptado de \textcite{Tran2022}.}

A convergência é rápida porque a jacobiana é bem condicionada e a solução varia pouco entre passos de tempo consecutivos, o que permite iniciar a iteração a partir do resultado anterior.
No caso particular do estado inicial, em que o campo de velocidades é nulo, a solução exata do sistema é $\boldsymbol{\lambda} = \mathbf{0}$ e o equilíbrio reduz-se à distribuição proporcional aos próprios pesos, $g_i^{eq} = 2\rho E\, W_i(T) / \sum_j W_j(T)$, que satisfaz a restrição de energia total.
Por ser executada em cada nó da malha e a cada passo de tempo, essa iteração impõe exigências específicas ao código, entre elas o critério de parada, a estimativa inicial e a proteção contra transbordamento numérico, tratadas na Seção~\ref{sec:implementacao-computacional-sequencial}.

%---------------------------------------------------------------------------------------

\subsection{Coeficientes de Transporte e Número de Prandtl}\label{subsec:lbm-transporte}

Os fatores de relaxação da \eqref{eq:ddf-f} e da \eqref{eq:ddf-g} vinculam-se aos coeficientes de transporte do fluido pelas relações
\begin{equation}\label{eq:coeficientes-transporte}
    \mu = \left( \frac{1}{\omega} - \frac{1}{2} \right) \rho T \Delta t, \qquad
    \kappa = c_p \left( \frac{1}{\omega_T} - \frac{1}{2} \right) \rho T \Delta t
\end{equation}
onde $\mu$ é a viscosidade dinâmica, $\kappa$ a condutividade térmica e $c_p$ o calor específico a pressão constante.
Invertendo essas expressões, obtêm-se os tempos de relaxação efetivamente calculados a cada nó,
\begin{equation}\label{eq:tempos-relaxacao}
    \tau = \frac{\mu}{\rho T} + \frac{1}{2}, \qquad
    \tau_T = \frac{\kappa}{\rho c_p T} + \frac{1}{2}
\end{equation}
que reproduzem a estrutura da \eqref{eq:viscosidade-tau}, agora com a temperatura local no lugar da velocidade do som constante.
O número de Prandtl resulta da razão entre ambos,
\begin{equation}\label{eq:prandtl}
    Pr = \frac{c_p \mu}{\kappa} = \frac{1/\omega - 1/2}{1/\omega_T - 1/2}
\end{equation}
e pode, portanto, assumir qualquer valor positivo.

Resta especificar a distribuição de quase equilíbrio $g_i^{*}$ que comparece na \eqref{eq:ddf-g}.
Sua função é corrigir o fluxo de calor de não equilíbrio, que de outro modo seria relaxado com a taxa da quantidade de movimento e não com a da energia.
Adota-se
\begin{equation}\label{eq:quase-equilibrio}
    g_i^{*} - g_i^{eq} = \frac{W_i(T)}{T} \left( \mathbf{c}_i \cdot \mathbf{L} \right), \qquad
    L_\alpha = 2 \sum_j \left( \mathbf{u} \cdot \mathbf{v}_j \right) v_{j\alpha} \left( f_j - f_j^{eq} \right)
\end{equation}
onde o vetor $\mathbf{L}$ é a contração da parcela de não equilíbrio do tensor de tensões com a velocidade macroscópica, ou seja, exatamente o trabalho das tensões viscosas que alimenta o fluxo de energia.
A correção anula-se quando o escoamento se encontra em equilíbrio local e quando $\omega_T = \omega$, situação em que a \eqref{eq:ddf-g} retorna à forma BGK simples e o número de Prandtl volta à unidade.

%---------------------------------------------------------------------------------------

\subsection{Estabilização por Relaxação Dependente do Não Equilíbrio Local}\label{subsec:lbm-sensor}

A positividade das populações e o desacoplamento térmico ainda não bastam para conter as oscilações que se formam junto às descontinuidades.
O mecanismo adicional proposto por \textcite{Tran2022} consiste em medir localmente o afastamento em relação ao equilíbrio e ajustar a relaxação em função dessa medida.
O indicador empregado é um número de Knudsen local, definido como o desvio relativo médio das populações:
\begin{equation}\label{eq:knudsen-local}
    Kn_{loc} = \frac{1}{Q} \sum_{i=0}^{Q-1} \frac{\left| f_i - f_i^{eq} \right|}{f_i^{eq}}
\end{equation}
Em regiões suaves do escoamento, $Kn_{loc}$ é pequeno, ao passo que nas proximidades de um choque ele se aproxima da unidade ou a supera.
O fator de relaxação efetivamente aplicado é então $\omega_L = \omega / \sigma(Kn_{loc})$, com $\sigma$ definido por partes conforme a \autoref{tab:lbm-sensor}, e limitado ao intervalo $1 \le \omega_L < 2$, faixa que assegura a estabilidade linear do esquema.

\begin{table}[H]
    \centering
    \caption{Fator de correção da relaxação em função do número de Knudsen local}
    \label{tab:lbm-sensor}
    \begin{tabular}{lcl}
        \toprule
        Faixa de $Kn_{loc}$        & $\sigma$ & Regime correspondente        \\
        \midrule
        $Kn_{loc} < 0{,}01$          & $1$      & Escoamento suave             \\
        $0{,}01 \le Kn_{loc} < 0{,}1$  & $1{,}05$   & Gradientes moderados         \\
        $0{,}1 \le Kn_{loc} < 1$     & $1{,}35$   & Vizinhança da descontinuidade \\
        $Kn_{loc} \ge 1$           & $\omega$ & Interior da frente de choque \\
        \bottomrule
    \end{tabular}
    \fonte{Adaptado de \textcite{Tran2022}.}
\end{table}

O comportamento pretendido é o de um mecanismo seletivo.
Onde a solução é suave, $\sigma = 1$ e a colisão permanece inalterada, de modo que a viscosidade física não é contaminada por dissipação artificial.
Onde há uma descontinuidade, $\sigma$ cresce, $\omega_L$ diminui e a relaxação torna-se menos rígida, o que equivale a uma viscosidade localmente maior e amortece as oscilações de Gibbs.
No limite $Kn_{loc} \ge 1$, a escolha $\sigma = \omega$ conduz a $\omega_L = 1$, ou seja, a colisão substitui integralmente as populações pelo equilíbrio, comportamento que corresponde ao limite invíscido e é o mais dissipativo admitido pelo esquema.

A combinação desses elementos responde às quatro limitações identificadas: a dupla distribuição fornece a temperatura como variável independente e a equação de estado dos gases ideais, as redes deslocadas reduzem o erro de invariância de Galileu, o segundo tempo de relaxação libera o número de Prandtl e o equilíbrio exponencial assegura a positividade.
O conjunto permite ao resolvedor operar em regimes de elevado número de Mach que estão fora do alcance do LBM clássico, ao custo de uma solução iterativa local em cada nó e de uma propagação interpolada.

%---------------------------------------------------------------------------------------

\subsection{Ciclo de Cálculo Resultante}\label{subsec:lbm-ciclo-compressivel}

O ciclo efetivamente implementado difere do ALGORITMO~\ref{alg:lbm-ciclo} em três aspectos, todos decorrentes dos elementos introduzidos nesta seção.
O cálculo dos equilíbrios deixa de ser uma avaliação polinomial direta e passa a incluir uma solução iterativa, a propagação se dá pela interpolação \eqref{eq:propagacao-idw} e a atualização das grandezas macroscópicas passa a fornecer também a temperatura e a pressão.
O ALGORITMO~\ref{alg:lbm-ciclo-compressivel} apresenta essa versão, cujas etapas correspondem uma a uma aos núcleos de cálculo descritos no Capítulo~\ref{ch:implementacao-computacional}.

\begin{algorithm}[H]
\caption{Ciclo de execução de um passo de tempo na formulação compressível}\label{alg:lbm-ciclo-compressivel}
\Entrada{Distribuições $f_i, g_i$ e campos macroscópicos $\rho, \mathbf{u}, T$ no instante $t$}
\Saida{Distribuições e campos macroscópicos no instante $t + \Delta t$}
\Para{cada nó $\mathbf{x}$ da malha}{
    Avaliar $W_i(T)$ pela \eqref{eq:pesos-temperatura} e $f_i^{eq}$ pela \eqref{eq:feq-produto}\;
    Obter $g_i^{eq}$ pelo ALGORITMO~\ref{alg:lbm-newton}\;
    Calcular $Kn_{loc}$ pela \eqref{eq:knudsen-local} e o fator $\omega_L$ correspondente\;
    Calcular $\tau$ e $\tau_T$ pela \eqref{eq:tempos-relaxacao} e a correção $g_i^{*} - g_i^{eq}$ pela \eqref{eq:quase-equilibrio}\;
    \tcp{Colisão}
    Relaxar $f_i$ e $g_i$ segundo a \eqref{eq:ddf-f} e a \eqref{eq:ddf-g}\;
}
\Para{cada nó $\mathbf{x}$ da malha}{
    \tcp{Propagação}
    Reunir $f_i$ e $g_i$ dos quatro nós de origem com os pesos da \eqref{eq:propagacao-idw}\;
}
\Para{cada nó $\mathbf{x}$ da malha}{
    \tcp{Grandezas macroscópicas}
    Calcular $\rho$, $\mathbf{u}$ e $T$ pela \eqref{eq:ddf-momentos} e pela \eqref{eq:temperatura-local}, e $p = \rho R T$\;
}
\Retorne{$f_i$, $g_i$, $\rho$, $\mathbf{u}$, $T$ e $p$}
\end{algorithm}
\fonte{O autor (2026), com base em \textcite{Tran2022} e \textcite{Rathnayake2024}.}

A separação do ciclo em três percursos da malha, adotada no ALGORITMO~\ref{alg:lbm-ciclo-compressivel} por clareza de exposição, decorre da mesma dependência de dados já identificada no modelo clássico, agora estendida à coleta interpolada, que exige a colisão concluída em toda a vizinhança antes de ter início.
O cálculo das grandezas macroscópicas, em contrapartida, é estritamente local e pode ser fundido à propagação, assim como a avaliação dos equilíbrios pode ser fundida à colisão, o que reduz o número de leituras e escritas na memória principal.
O número de percursos efetivamente realizados, bem como a quantificação do custo por nó, do volume de dados movimentado e da intensidade aritmética resultante, são tratados na Seção~\ref{sec:custo-computacional}, no âmbito das decisões de implementação a que essas grandezas dizem respeito.
Cabe registrar que a formulação compressível preserva as duas propriedades que tornam o LBM adequado ao paralelismo massivo, uma vez que a colisão permanece inteiramente local, sem qualquer dependência entre nós, e a propagação alcança apenas os vizinhos imediatos, ainda que por meio de pesos de interpolação.

%---------------------------------------------------------------------------------------

\section{Tratamento de Condições de Contorno}\label{sec:lbm-contorno-hipersonico}

O tratamento de condições de contorno no LBM exige a determinação das funções de distribuição que entram no domínio a partir das fronteiras, de modo que as leis de conservação e o comportamento hidrodinâmico sejam preservados.
Entre as técnicas mais comuns está a reflexão de partículas (\textit{bounce-back}), utilizada para impor a condição de não escorregamento em paredes sólidas, na qual as distribuições que atingem a parede são refletidas de volta para o domínio, sendo essa técnica valorizada por sua simplicidade e por sua capacidade de tratar geometrias complexas sem perda de massa \parencite{He1997}.
Destacam-se também as condições de contorno de equilíbrio e de não equilíbrio, nas quais as distribuições na fronteira são impostas combinando os valores de equilíbrio correspondentes às variáveis macroscópicas desejadas com correções de não equilíbrio, sendo o método proposto por \textcite{Zou1997} amplamente utilizado para prescrever perfis de velocidade ou de pressão com precisão de segunda ordem.
Por fim, a reflexão antiespecular (\textit{anti-specular bounce-back}) é empregada em fronteiras de saída para minimizar reflexões de ondas de pressão e garantir a saída suave do fluido do domínio computacional.

No caso do tubo de choque simulado neste trabalho, a geometria dispensa paredes sólidas e o tratamento reduz-se a duas condições.
Nas extremidades do tubo, aplica-se a condição de gradiente nulo, implementada pela limitação do índice do nó de origem ao intervalo válido da malha, o que equivale a estender indefinidamente o estado da célula de borda e é compatível com o fato de as ondas não alcançarem as extremidades no intervalo de tempo simulado.
Na direção transversal, adota-se periodicidade, o que reproduz um escoamento invariante nessa direção e torna cada linha da malha comparável à solução analítica unidimensional.
Ambas as condições são incorporadas às próprias tabelas de propagação, de modo que a etapa de propagação permanece livre de desvios condicionais, propriedade relevante para as implementações em aceleradores de \textit{hardware}.

A condição de gradiente nulo nas extremidades e a periodicidade na direção transversal delimitam, assim, um problema em que o tratamento das fronteiras não interfere na dinâmica das ondas durante o intervalo simulado, o que isola a comparação entre tecnologias de paralelização de qualquer efeito associado ao contorno.
